% =============================================================================
% Chuong 3: Cong trinh lien quan
% =============================================================================
% Cach dung tren Overleaf:
%   1. Bao dam da tai macros.tex, refs.bib va chapters/ch3.tex vao du an.
%   2. Trong main.tex:
%        % =============================================================================
% Macro dung chung cho luan van (theo docs/drafts/style-guide-luan-van.md §6)
% =============================================================================
% Them cac goi sau vao main.tex tren Overleaf (neu template chua co):
%   \usepackage[utf8]{inputenc}   % (bo qua neu dung XeLaTeX / LuaLaTeX)
%   \usepackage[vietnamese]{babel} % hoac vntex / polyglossia
%   \usepackage{amsmath,amssymb}
%   \usepackage{booktabs}
%   \usepackage{tabularx}
%   \usepackage{array}
%   \usepackage{graphicx}
%   \usepackage{xcolor}
%   \usepackage{icomma}           % dau phay thap phan trong moi truong toan
%   \usepackage{hyperref}
% =============================================================================

\newcommand{\CTAV}{CTAV}
\newcommand{\code}[1]{\texttt{#1}}
\newcommand{\splitname}[1]{\texttt{#1}}
\newcommand{\CHO}[2]{\textcolor{red}{[CHỜ #1: #2]}}

% Cot tu dong xuong dong can le trai cho tabularx
\newcolumntype{Y}{>{\raggedright\arraybackslash}X}

%        ...
%        % =============================================================================
% Chuong 3: Cong trinh lien quan
% =============================================================================
% Cach dung tren Overleaf:
%   1. Bao dam da tai macros.tex, refs.bib va chapters/ch3.tex vao du an.
%   2. Trong main.tex:
%        % =============================================================================
% Macro dung chung cho luan van (theo docs/drafts/style-guide-luan-van.md §6)
% =============================================================================
% Them cac goi sau vao main.tex tren Overleaf (neu template chua co):
%   \usepackage[utf8]{inputenc}   % (bo qua neu dung XeLaTeX / LuaLaTeX)
%   \usepackage[vietnamese]{babel} % hoac vntex / polyglossia
%   \usepackage{amsmath,amssymb}
%   \usepackage{booktabs}
%   \usepackage{tabularx}
%   \usepackage{array}
%   \usepackage{graphicx}
%   \usepackage{xcolor}
%   \usepackage{icomma}           % dau phay thap phan trong moi truong toan
%   \usepackage{hyperref}
% =============================================================================

\newcommand{\CTAV}{CTAV}
\newcommand{\code}[1]{\texttt{#1}}
\newcommand{\splitname}[1]{\texttt{#1}}
\newcommand{\CHO}[2]{\textcolor{red}{[CHỜ #1: #2]}}

% Cot tu dong xuong dong can le trai cho tabularx
\newcolumntype{Y}{>{\raggedright\arraybackslash}X}

%        ...
%        % =============================================================================
% Chuong 3: Cong trinh lien quan
% =============================================================================
% Cach dung tren Overleaf:
%   1. Bao dam da tai macros.tex, refs.bib va chapters/ch3.tex vao du an.
%   2. Trong main.tex:
%        % =============================================================================
% Macro dung chung cho luan van (theo docs/drafts/style-guide-luan-van.md §6)
% =============================================================================
% Them cac goi sau vao main.tex tren Overleaf (neu template chua co):
%   \usepackage[utf8]{inputenc}   % (bo qua neu dung XeLaTeX / LuaLaTeX)
%   \usepackage[vietnamese]{babel} % hoac vntex / polyglossia
%   \usepackage{amsmath,amssymb}
%   \usepackage{booktabs}
%   \usepackage{tabularx}
%   \usepackage{array}
%   \usepackage{graphicx}
%   \usepackage{xcolor}
%   \usepackage{icomma}           % dau phay thap phan trong moi truong toan
%   \usepackage{hyperref}
% =============================================================================

\newcommand{\CTAV}{CTAV}
\newcommand{\code}[1]{\texttt{#1}}
\newcommand{\splitname}[1]{\texttt{#1}}
\newcommand{\CHO}[2]{\textcolor{red}{[CHỜ #1: #2]}}

% Cot tu dong xuong dong can le trai cho tabularx
\newcolumntype{Y}{>{\raggedright\arraybackslash}X}

%        ...
%        % =============================================================================
% Chuong 3: Cong trinh lien quan
% =============================================================================
% Cach dung tren Overleaf:
%   1. Bao dam da tai macros.tex, refs.bib va chapters/ch3.tex vao du an.
%   2. Trong main.tex:
%        \input{macros}
%        ...
%        \include{chapters/ch3}
%        ...
%        \bibliographystyle{IEEEtran}
%        \bibliography{refs}
%
% Cac nhan lien chuong duoc tham chieu tu Chuong 3:
%   Chuong: ch:mo-dau, ch:nen-tang, ch:de-xuat, ch:thuc-nghiem, ch:platform
%   Chuong 2: sec:ch2-du-lieu, sec:ch2-truy-xuat, sec:ch2-judge, sec:ch2-rca,
%             sec:ch2-dang-ky
%   Chuong 4: sec:ch4-ctav, sec:ch4-c1, sec:ch4-policy, sec:ch4-c4, sec:ch4-rca
%   Chuong 5: sec:ch5-hypotheses, sec:ch5-c5, sec:ch5-budget
% =============================================================================

\chapter{Công trình liên quan}
\label{ch:cong-trinh-lien-quan}

Chương này hệ thống hóa các hướng nghiên cứu và các hệ thống vận hành thực tế liên quan trực tiếp tới bài toán thẩm định cảnh báo có nhận biết topology và phân tích nguyên nhân gốc trên kiến trúc microservice. Nội dung được tổ chức xoay quanh hai câu hỏi: cộng đồng học thuật và công nghiệp đã giải quyết những khía cạnh nào của chuỗi xử lý từ phát hiện bất thường đến định vị nguyên nhân gốc, và còn những khoảng trống nào mà đề xuất của đề tài nhắm tới. Mục~\ref{sec:ch3-nghien-cuu} khảo sát bảy nhánh nghiên cứu học thuật từ phát hiện bất thường chuỗi thời gian, ứng dụng mô hình ngôn ngữ lớn trong phân tích log và nguyên nhân gốc, kiến trúc giám sát lai, phân tích nguyên nhân gốc dựa trên đồ thị, truy xuất dựa trên đồ thị, hiệu chỉnh độ tin cậy của \emph{LLM-as-a-Judge}, đến gom nhóm và phân loại cảnh báo. Mục~\ref{sec:ch3-he-thong} đối chiếu các nền tảng AIOps thương mại và công cụ mã nguồn mở hiện hành dựa trên tài liệu chính thức. Cuối cùng, Mục~\ref{sec:ch3-ket-chuong} tổng hợp các hướng tiếp cận trong Bảng~\ref{tab:ch3-so-sanh}, rút ra bốn khoảng trống nghiên cứu và định vị rõ phạm vi đóng góp của \CTAV{} cùng tác tử phân tích nguyên nhân gốc.

% -----------------------------------------------------------------------------
\section{Các nghiên cứu liên quan}
\label{sec:ch3-nghien-cuu}

Các công trình liên quan đến đề tài nằm ở giao điểm giữa giám sát chuỗi thời gian, phân tích dữ liệu quan sát đa phương thức, truy xuất bằng chứng theo đồ thị và hiệu chỉnh mô hình ngôn ngữ lớn đóng vai trò bộ thẩm định. Bảy tiểu mục dưới đây phân tích từng nhánh theo cùng một trục: cách tiếp cận của công trình, hạn chế đặt trong bài toán của đề tài và điểm khác biệt của đề tài.

\subsection{Phát hiện bất thường trên chuỗi thời gian metric}
\label{sec:ch3-tsad}

Hướng phát hiện bất thường trên chuỗi thời gian không giám sát hình thành nền tảng cho tầng lọc cảnh báo đầu vào. Liu và cộng sự~\cite{liu2008iforest} đề xuất Isolation Forest, một phương pháp dựa trên tập cây phân tách ngẫu nhiên tận dụng nhận xét rằng điểm bất thường ít và khác biệt nên cần ít phép chia để cô lập hơn điểm bình thường; mô hình chạy nhanh, không cần nhãn và tiêu thụ ít tài nguyên. Mở rộng sang dữ liệu luồng có phân bố biến đổi theo thời gian, Guha và cộng sự~\cite{guha2016rrcf} xây dựng Robust Random Cut Forest (RRCF) cho phép cập nhật cấu trúc cây trực tuyến khi điểm mới đến và đo độ bất thường bằng mức thay đổi độ phức tạp cây khi chèn hoặc xóa điểm. Hạn chế chung của nhóm cây phân tách là chỉ nhìn từng vector metric độc lập theo cửa sổ, nên không phân biệt được dao động tải bình thường với sự cố thực sự nếu không đối chiếu ngữ cảnh văn bản từ log hoặc quan hệ gọi giữa các service.

Để nắm bắt tương quan đa biến và phụ thuộc giữa các cảm biến hay metric, các công trình học sâu dựa trên đồ thị đã được đề xuất. Zhao và cộng sự~\cite{zhao2020mtadgat} giới thiệu MTAD-GAT với hai lớp chú ý đồ thị song song: một lớp học tương quan giữa các đặc trưng tại cùng thời điểm và một lớp học phụ thuộc dọc trục thời gian, kết hợp đồng thời mục tiêu dự báo và tái tạo để tính điểm bất thường. Deng và Hooi~\cite{deng2021gdn} đề xuất Graph Deviation Network (GDN), tự học đồ thị quan hệ giữa các cảm biến qua vector nhúng và dùng mạng chú ý đồ thị để dự báo giá trị tiếp theo, từ đó giải thích bất thường qua độ lệch của các cạnh dự báo. Tuy đạt độ chính xác cao hơn trên nhiều tập chuỗi thời gian đa biến, nhóm học sâu đồ thị đòi hỏi huấn luyện lại khi topology hoặc tập metric thay đổi, cần tài nguyên tính toán lớn hơn đáng kể và vẫn chỉ khai thác tín hiệu số mà bỏ qua nội dung log hay trace vốn chứa thông tin tường minh về lỗi ứng dụng.

Khác với các công trình tập trung thay thế bộ phát hiện metric bằng kiến trúc học sâu ngày càng phức tạp, đề tài này giữ tầng phát hiện nhẹ bằng Isolation Forest để bảo đảm tốc độ và khả năng chạy lại tức thì trên mọi cấu hình, rồi chuyển trọng tâm sang tầng thẩm định bằng bằng chứng đa phương thức phía sau. Cách phân công này cho phép tận dụng tốc độ lọc sơ bộ của detector nhẹ trong khi vẫn kiểm soát được báo động giả nhờ đối chiếu log và đồ thị phụ thuộc service.

\subsection{Mô hình ngôn ngữ lớn trong phân tích log và nguyên nhân gốc}
\label{sec:ch3-llm-log-rca}

Sự xuất hiện của các mô hình ngôn ngữ lớn (LLM) mở ra hướng xử lý trực tiếp dữ liệu log phi cấu trúc và suy luận sự cố bằng ngôn ngữ tự nhiên. Ở bài toán phân tách mẫu log (\emph{log parsing}), Zhong và cộng sự~\cite{zhong2024logparser} đề xuất LogParser-LLM, kết hợp cây tiền tố với truy vấn LLM chọn lọc để trích xuất mẫu log mà không cần huấn luyện lại mô hình khi định dạng log thay đổi. Ở bài toán phát hiện bất thường trên chuỗi log, Guan và cộng sự~\cite{guan2024logllm} giới thiệu LogLLM, phối hợp bộ mã hóa ngữ nghĩa dạng BERT cho từng dòng log với bộ giải mã LLM nhằm mô hình hóa chuỗi log mà không cần bước tách mẫu thủ công. Hai công trình này giải quyết tốt khâu biểu diễn và phân tích log đơn lẻ, nhưng chỉ làm việc trên tín hiệu log tách rời mà không tích hợp luồng cảnh báo từ metric cũng như cấu trúc phụ thuộc giữa các service.

Ở bài toán phân tích nguyên nhân gốc cho sự cố đám mây, nhiều công trình sử dụng LLM trực tiếp trên hồ sơ sự cố hoặc dưới dạng tác tử gọi công cụ. Ahmed và cộng sự~\cite{ahmed2023recommending} thực hiện đánh giá quy mô lớn trên hơn 40.000 sự cố vận hành thực tế tại Microsoft để khuyến nghị nguyên nhân gốc và bước khắc phục bằng mô hình dòng GPT, cho thấy LLM có ích trong việc tóm tắt tri thức sự cố nhưng nhạy cảm với chất lượng mô tả đầu vào. Chen và cộng sự~\cite{chen2024rcacopilot} xây dựng RCACopilot, một hệ thống định tuyến sự cố vào quy trình thu thập dữ liệu chẩn đoán tương ứng rồi để LLM dự đoán nhóm nguyên nhân gốc kèm giải thích bằng ngôn ngữ tự nhiên. Tiến thêm bước tự chủ, Wang và cộng sự~\cite{wang2024rcagent} đề xuất RCAgent sử dụng tác tử LLM có hỗ trợ công cụ truy vấn log và tổng hợp tự nhất quán (\emph{self-consistency}) trong môi trường đám mây, trong khi Roy và cộng sự~\cite{roy2024exploring} khảo sát tác tử theo vòng lặp ReAct trang bị công cụ truy xuất log, metric và cơ sở dữ liệu lịch sử để điều tra sự cố.

Điểm hạn chế chung của nhóm công trình LLM cho RCA nêu trên là phần lớn giả định sự cố đầu vào đã được xác nhận hoặc đã mở phiếu sự cố sẵn, nên không đánh giá trực tiếp bài toán thẩm định luồng cảnh báo thô để chặn báo động giả trước khi gọi tác tử chẩn đoán tốn kém. Ngoài ra, các hệ thống đánh giá trên dữ liệu sự cố nội bộ doanh nghiệp khó tái lập độc lập bên ngoài. Đề tài này đặt tầng thẩm định cảnh báo có kiểm soát ngân sách ngay trước tác tử RCA và thực nghiệm hoàn toàn trên bộ dữ liệu công khai có thể tái lập.

\subsection{Kiến trúc giám sát lai và LLM-as-a-Judge}
\label{sec:ch3-hybrid-judge}

Để dung hòa chi phí gọi LLM với yêu cầu giám sát liên tục, nhiều công trình kết hợp mô hình nhẹ hoặc cơ chế truy xuất tăng cường sinh (\emph{Retrieval-Augmented Generation}, RAG)~\cite{lewis2020rag} với bộ phân xử LLM ở tầng sau. Về nguyên lý phân tầng chi phí, nền tảng lý thuyết bắt nguồn từ cơ chế đánh đổi sai số và quyền từ chối (\emph{reject option}) của Chow~\cite{chow1970reject} cùng mô hình chuỗi tầng gọi LLM tiết kiệm chi phí FrugalGPT của Chen và cộng sự~\cite{chen2024frugalgpt}, trong đó các mẫu rõ ràng được xử lý bởi mô hình rẻ hơn và chỉ vùng bất định mới chuyển lên mô hình đắt hơn.

Trong miền phát hiện bất thường, một loạt công trình gần đây đưa kiến trúc lai vào thực nghiệm. Russell-Gilbert và cộng sự~\cite{russellgilbert2025raad} đề xuất RAAD-LLM, tích hợp RAG vào LLM để thích nghi phát hiện bất thường chuỗi thời gian khi bối cảnh hệ thống thay đổi mà không cần tinh chỉnh trọng số. Gu và cộng sự~\cite{gu2025argos} giới thiệu ARGOS, một khuôn khổ tác tử cho phép LLM tự sinh và tinh chỉnh luật phát hiện bất thường trên chuỗi thời gian nhằm giảm phụ thuộc vào truy vấn trực tiếp trên từng mẫu. Devireddy và Huang~\cite{devireddy2025calm} đề xuất CALM, điều phối liên tục giữa phát hiện thời gian thực và cơ chế can thiệp thích nghi của LLM. Ở góc độ đánh giá năng lực thẩm định, Sivalingam và cộng sự~\cite{sivalingam2026tsquery} xây dựng TSQueryBench để kiểm tra khả năng dùng \emph{LLM-as-a-Judge} đánh giá lời giải thích bất thường chuỗi thời gian.

Hạn chế chung của các kiến trúc lai này khi áp vào microservice là chúng chủ yếu vận hành trên chuỗi thời gian đơn lẻ hoặc truy xuất ngữ cảnh văn bản phẳng, chưa kết nối tường minh với đồ thị phụ thuộc giữa các service sinh từ trace. Quan trọng hơn, các nghiên cứu này chưa khảo sát hệ quả ngược chiều khi mở rộng bằng chứng cho bộ thẩm định LLM: việc cấp thêm ngữ cảnh từ láng giềng có thể làm tăng tỷ lệ báo động giả do thiên lệch xác nhận. Đề tài này tập trung giải quyết trực tiếp bài toán kiểm soát ngân sách và kìm báo động giả trong kiến trúc lai detector nhẹ kết hợp \emph{LLM-as-a-Judge} cho hệ microservice.

\subsection{Phân tích nguyên nhân gốc dựa trên topology, đồ thị và suy luận nhân quả}
\label{sec:ch3-rca-graph}

Định vị nguyên nhân gốc dựa trên đồ thị phụ thuộc và suy luận nhân quả là nhánh phát triển lâu đời nhất trong giám sát hệ phân tán. Kim và cộng sự~\cite{kim2013monitorrank} đề xuất MonitorRank, thực hiện bước đi ngẫu nhiên (\emph{random walk}) có trọng số tương quan metric trên đồ thị gọi dịch vụ để xếp hạng service nghi vấn trong kiến trúc hướng dịch vụ. Chen và cộng sự~\cite{chen2014causeinfer} xây dựng CauseInfer với đồ thị nhân quả hai tầng (tầng dịch vụ và tầng metric) bằng thuật toán kiểm định độc lập có điều kiện PC để lần ngược nút gốc gây suy giảm hiệu năng. Trên kiến trúc microservice, Wu và cộng sự~\cite{wu2020microrca} giới thiệu MicroRCA, gắn dữ liệu sử dụng tài nguyên máy chủ vào đồ thị dịch vụ bất thường và áp dụng Personalized PageRank để định vị service gây lỗi hiệu năng.

Các công trình gần đây tiếp cận bài toán bằng mô hình nhân quả chặt chẽ hơn và dữ liệu đa phương thức. Li và cộng sự~\cite{li2022circa} đề xuất CIRCA, mô hình hóa lỗi như một can thiệp (\emph{intervention}) trên mạng Bayes nhân quả xây từ kiến trúc hệ thống và dùng điều kiện thay đổi phân bố có điều kiện để chỉ ra chỉ số gốc. Ikram và cộng sự~\cite{ikram2022rcd} đề xuất RCD, xem sự cố như can thiệp cục bộ và khám phá đồ thị nhân quả theo phân cấp để tránh chi phí tính toán lớn của thuật toán PC trên toàn bộ hệ thống. Kết hợp cả ba trụ cột dữ liệu quan sát, Yu và cộng sự~\cite{yu2023nezha} đề xuất Nezha, chuyển đổi metric, log và trace thành biểu diễn sự kiện đồng nhất và so sánh mẫu phân kỳ giữa trạng thái bình thường với trạng thái lỗi để định vị nguyên nhân gốc ở mức chi tiết, đồng thời công bố bộ dữ liệu microservice đa phương thức làm chuẩn đánh giá. Đáng lưu ý về tính bền vững của chuỗi xử lý, Pham và cộng sự~\cite{pham2024baro} đề xuất BARO bằng phát hiện điểm đổi Bayes đa biến kết hợp kiểm định phi tham số, đồng thời chỉ ra rằng sai số ở khâu phát hiện bất thường đầu vào có thể lan sang làm giảm mạnh độ chính xác định vị gốc; tiếp đó nhóm tác giả công bố bộ chuẩn RCAEval~\cite{pham2025rcaeval} tổng hợp 735 ca lỗi đa phương thức trên nhiều hệ microservice mã nguồn mở để thống nhất việc so sánh.

Hạn chế của nhóm phương pháp đồ thị và nhân quả thuần túy gồm hai điểm. Thứ nhất, phần lớn mô hình chỉ xuất ra danh sách xếp hạng service hoặc metric gốc dưới dạng chỉ số toán học, thiếu báo cáo tường minh giải thích cơ chế lan truyền bằng chứng từ log và trace cho kỹ sư vận hành. Thứ hai, trừ những phân tích gần đây như BARO, đa số nghiên cứu RCA giả định cửa sổ thời gian sự cố đã được xác định sạch từ trước, tách rời khỏi bài toán lọc cảnh báo giả trực tuyến. Đề tài này nối trực tiếp hai khâu: dùng \CTAV{} để làm sạch luồng cảnh báo trước khi gom sự cố, rồi dùng tác tử RCA lần ngược chuỗi lan truyền trên đồ thị phụ thuộc và sinh báo cáo có cấu trúc kèm trích dẫn bằng chứng.

\subsection{Truy xuất dựa trên đồ thị và GraphRAG cho vận hành hệ thống}
\label{sec:ch3-graphrag}

Trong lĩnh vực xử lý ngôn ngữ tự nhiên, Edge và cộng sự~\cite{edge2024graphrag} đề xuất GraphRAG nhằm khắc phục hạn chế của truy xuất văn bản phẳng khi cần tổng hợp thông tin phân tán trên nhiều tài liệu. Phương pháp xây dựng đồ thị tri thức thực thể--quan hệ từ kho văn bản, phân cụm cộng đồng đồ thị và sinh tóm tắt phân cấp để trả lời các truy vấn mang tính toàn cục. Tuy mở ra hướng truy xuất có cấu trúc, GraphRAG thiết kế cho kho tài liệu tĩnh với đồ thị trích xuất bằng LLM, nên vừa tốn chi phí lập chỉ mục vừa khác về bản chất so với đồ thị phụ thuộc động của hệ microservice được đo trực tiếp từ quan hệ gọi \code{SpanID}--\code{ParentID} trong trace.

Khi áp dụng tư tưởng truy xuất theo cấu trúc đồ thị vào bài toán thẩm định cảnh báo microservice, một nghiên cứu tiền đề về truy xuất bằng chứng có nhận biết topology~\cite{author2026icaetm} đã đề xuất kết hợp điểm từ khóa BM25 với độ gần về thời gian và khoảng cách bước nhảy trên đồ thị phụ thuộc service để cung cấp thêm bằng chứng từ các service láng giềng cho bộ thẩm định LLM (chi tiết công thức hợp nhất đã trình bày ở Mục~\ref{sec:ch2-truy-xuat}). Kết quả định tính của hướng tiếp cận này cho thấy một đánh đổi rõ rệt: việc mở rộng ngữ cảnh sang các service có liên kết gọi giúp cải thiện khả năng nhận diện các sự cố lan truyền mà bằng chứng cục bộ không đủ, nhưng khi đưa trực tiếp vào \emph{LLM-as-a-Judge} mà không có cơ chế kiểm soát chính sách, số báo động giả tăng lên mạnh do bộ thẩm định dễ diễn giải nhầm log nền từ láng giềng thành dấu hiệu bất thường. Đề tài này xuất phát trực tiếp từ đánh đổi đó để xây dựng cơ chế hiệu chỉnh theo nguồn gốc bằng chứng và kiểm định nghiêm ngặt hiệu quả kìm báo động giả trên dữ liệu ngoài tập tinh chỉnh.

\subsection{Hiệu chỉnh độ tin cậy, thiên lệch của LLM-as-a-Judge và gắn bằng chứng}
\label{sec:ch3-calibration}

Dù \emph{LLM-as-a-Judge} đạt mức đồng thuận cao với người đánh giá trên nhiều tác vụ tổng quát~\cite{zheng2023judging}, hàng loạt nghiên cứu đã chứng minh mô hình ngôn ngữ mắc nhiều thiên lệch có hệ thống và thường tự tin thái quá. Wang và cộng sự~\cite{wang2024fair} chỉ ra rằng bộ thẩm định LLM chịu thiên lệch vị trí nghiêm trọng khi so sánh cặp câu trả lời, dẫn đến thứ hạng dễ bị thao túng chỉ bằng cách đảo thứ tự trình bày, và đề xuất khung FairEval hiệu chỉnh bằng cách buộc sinh bằng chứng trước khi chấm điểm kết hợp trung bình hóa vị trí. Sharma và cộng sự~\cite{sharma2024sycophancy} phân tích hiện tượng chiều lòng (\emph{sycophancy}), cho thấy mô hình đã qua tinh chỉnh theo phản hồi con người có xu hướng đồng thuận với giả định ngầm trong câu hỏi ngay cả khi giả định đó sai — một cơ chế tương đồng trực tiếp với thiên lệch xác nhận khi bộ thẩm định nhận thông báo rằng ``cảnh báo đang kích hoạt''.

Về hiệu chỉnh xác suất và ước lượng độ bất định, Guo và cộng sự~\cite{guo2017calibration} chứng minh các mạng nơ-ron hiện đại tuy chính xác nhưng hiệu chỉnh kém và đề xuất phương pháp hiệu chỉnh hậu xử lý bằng tỉ lệ nhiệt độ (\emph{temperature scaling}). Trên LLM, Kadavath và cộng sự~\cite{kadavath2022mostly} khảo sát khả năng mô hình tự đánh giá xác suất đúng của phát biểu do chính nó sinh ra, cho thấy tín hiệu nội tại có thể hiệu chỉnh được trong câu hỏi trắc nghiệm nhưng khó khái quát khi chuyển phân bố tác vụ. Gần hơn với bài toán sử dụng API thương mại không truy cập được xác suất token thô, Xiong và cộng sự~\cite{xiong2024uncertainty} đánh giá thực nghiệm độ tin cậy bằng lời (\emph{verbalized confidence}) trên nhiều dòng LLM và chỉ ra mô hình thường gán điểm tin cậy quá cao, tập trung dày đặc ở vùng sát $1{,}0$. Song song đó, để hạn chế suy diễn vô căn cứ trong văn bản sinh ra, Gao và cộng sự~\cite{gao2023alce} đề xuất chuẩn đánh giá ALCE đo đồng thời tính đúng đắn và chất lượng trích dẫn nguồn bằng chứng trong câu trả lời của LLM.

Khoảng trống còn lại trong nhánh nghiên cứu này là các kỹ thuật hiệu chỉnh hiện có chủ yếu thiết kế cho bài toán chấm điểm cặp câu trả lời, hỏi--đáp mở hoặc sinh văn bản có trích dẫn tổng quát. Chưa có cơ chế nào hiệu chỉnh trực tiếp quyết định của \emph{LLM-as-a-Judge} trong bài toán thẩm định cảnh báo microservice theo \emph{nguồn gốc bằng chứng} (nội bộ hay láng giềng trên đồ thị phụ thuộc) và theo \emph{mức độ bất thường nền đo trên tập phát triển}. Đề tài này cụ thể hóa ý tưởng hiệu chỉnh bằng chính sách neo bằng chứng và kiểm định tách biệt từng thành phần thiết kế (ngưỡng bất đối xứng theo nguồn bằng chứng, neo ngưỡng nền trên tập phát triển và định tuyến hai tầng) theo quy trình đăng ký giả thuyết trước ở Mục~\ref{sec:ch2-dang-ky}.

\subsection{Gom nhóm, phân loại và cứu cảnh báo trong vận hành}
\label{sec:ch3-alert-triage}

Trong vận hành hệ thống trực tuyến quy mô lớn, xử lý bão cảnh báo và định tuyến sự cố là bài toán sống còn để giảm tải cho kỹ sư trực. Zhao và cộng sự~\cite{zhao2020alertstorm} khảo sát hiện tượng bão cảnh báo trong hệ dịch vụ trực tuyến công nghiệp và đề xuất quy trình hai bước gồm phát hiện thời điểm bão và gom nhóm các cảnh báo tương quan về thời gian--ngữ nghĩa thành một cụm sự cố đại diện. Ở khâu định tuyến sự cố, Chen và cộng sự~\cite{chen2019empirical,chen2019continuous} phân tích nguyên nhân định tuyến nhầm sự cố tại các dịch vụ đám mây lớn và xây dựng phương pháp phân loại liên tục tận dụng thông tin cập nhật trong quá trình thảo luận của kỹ sư. Gần đây hơn, Kuang và cộng sự~\cite{kuang2024knowledge} đề xuất phương pháp gom cảnh báo có nhận biết tri thức cho hệ đám mây quy mô lớn, kết hợp đặc trưng văn bản với tri thức vận hành và topology để gộp cảnh báo cùng gốc tốt hơn các luật tĩnh.

Các công trình gom nhóm và phân loại này giải quyết tốt khâu \emph{sau khi cảnh báo đã phát sinh}, giúp nén số lượng thông báo gửi tới người trực, nhưng thường không thẩm định lại bản thân cảnh báo đầu vào có phải là báo động giả do bộ phát hiện metric gây ra hay không. Mặt khác, một câu hỏi ít được kiểm chứng nghiêm ngặt là liệu các cửa sổ bị bộ phát hiện chính bỏ sót có thể được ``cứu'' lại bằng điểm phụ trợ từ tín hiệu log/trace (cơ chế \code{Rescue Score}) mà không làm vỡ tỷ lệ báo động giả hay không. Đề tài này tách rõ hai vấn đề: sử dụng gom nhóm theo cửa sổ và topology ở đầu vào tác tử RCA để hợp nhất các cảnh báo đã được \CTAV{} thẩm định thành sự cố, đồng thời đánh giá cơ chế cứu cảnh báo như một phân tích loại trừ độc lập (Mục~\ref{sec:ch5-c5}) thay vì mặc định đưa vào hệ thống chính.

% -----------------------------------------------------------------------------
\section{Các hệ thống và công cụ hiện có}
\label{sec:ch3-he-thong}

Bên cạnh các công trình học thuật, nhiều nền tảng AIOps thương mại và công cụ mã nguồn mở đã được triển khai rộng rãi trong thực tế vận hành. Mục này mô tả tính năng cốt lõi của các hệ thống tiêu biểu dựa trên tài liệu chính thức của nhà phát triển, sau đó chỉ ra điểm khác biệt về mục tiêu và giới hạn khi đặt trong bối cảnh nghiên cứu của đề tài.

\subsection{Nền tảng AIOps thương mại}
\label{sec:ch3-thuong-mai}

\textbf{Datadog (Event Management và Bits AI).} Theo tài liệu chính thức của Datadog~\cite{datadog2026docs}, phân hệ Event Management thu nạp tập trung cảnh báo từ các bộ giám sát nội bộ cùng nguồn bên thứ ba, sau đó áp dụng tương quan cảnh báo tự động dựa trên topology dịch vụ, nhãn chia sẻ và mẫu thời gian để gom các cảnh báo lan truyền vào chung một sự vụ, đồng thời làm giàu ngữ cảnh về đơn vị sở hữu dịch vụ. Ở lớp điều tra sự cố, tác tử Bits AI (Bits Investigation) được kích hoạt tự động khi có cảnh báo nghiêm trọng, thực hiện kiểm định song song nhiều giả thuyết trên dữ liệu metric, log, trace và thay đổi cấu hình để lần theo chuỗi nhân quả và đề xuất hướng xử lý. Tuy cung cấp quy trình khép kín từ gom cảnh báo đến điều tra bằng AI, kiến trúc bên trong của cơ chế tương quan và điều tra đều đóng kín; nhà cung cấp không công bố phương pháp hiệu chỉnh độ tin cậy của bộ suy luận LLM theo nguồn bằng chứng hay đánh giá tách biệt tác động của topology lên tỷ lệ báo động giả trên bộ chuẩn công khai.

\textbf{Dynatrace (Davis AI và Smartscape).} Tài liệu kỹ thuật của Dynatrace~\cite{dynatrace2026docs} mô tả kiến trúc dựa trên hai trụ cột là bản đồ phụ thuộc thời gian thực Smartscape và động cơ phân tích nhân quả Davis AI. Smartscape tự động lập đồ thị phụ thuộc đầy đủ từ ứng dụng, service, tiến trình đến máy chủ và hạ tầng Kubernetes. Khi phát sinh bất thường ở mức chỉ số, độ trễ hoặc tỷ lệ lỗi, Davis AI gom các sự kiện liên quan có cùng nguyên nhân gốc thành một thực thể \emph{Problem} duy nhất và thực hiện phân tích cây lỗi tất định (\emph{deterministic fault-tree analysis}) trên đồ thị Smartscape để tách biệt thành phần gốc khỏi các thành phần chịu ảnh hưởng, hiển thị qua biểu đồ \emph{Visual Resolution Path}. Ưu điểm của Dynatrace là tính nhất quán cao nhờ mô hình nhân quả dựa trên agent thu thập sâu (\emph{OneAgent}), nhưng hệ thống gắn chặt với hạ tầng thu thập độc quyền, chi phí bản quyền lớn và không nhằm giải bài toán hiệu chỉnh quyết định của \emph{LLM-as-a-Judge} trên dữ liệu log phi cấu trúc lẫn ngân sách truy vấn LLM giới hạn.

\textbf{PagerDuty AIOps.} Theo tài liệu hỗ trợ chính thức của PagerDuty~\cite{pagerduty2026docs}, bộ công cụ PagerDuty AIOps tập trung vào ba chức năng: \emph{Intelligent Alert Grouping} dùng học máy nhóm các cảnh báo có liên quan theo mẫu lịch sử để giảm nhiễu thông báo; \emph{Probable Origin} phân tích chuỗi sự cố lan truyền trong lịch sử sáu tháng gần nhất để gợi ý tối đa ba dịch vụ hoặc sự cố có khả năng là điểm khởi phát cao nhất; và \emph{Event Orchestration} cung cấp bộ máy luật cho phép lọc, làm giàu và định tuyến sự kiện tự động ở mức toàn cục hoặc từng dịch vụ. Cách tiếp cận của PagerDuty giải quyết hiệu quả bài toán quản trị phản ứng sự cố ở tầng thông báo, nhưng phụ thuộc nhiều vào lịch sử sự cố tích lũy theo tổ chức, không khai thác đồ thị gọi chi tiết từng cửa sổ trace để thẩm định cảnh báo từ tín hiệu thô và không công bố cơ chế kiểm soát báo động giả cho tác tử thẩm định LLM.

\subsection{Công cụ mã nguồn mở}
\label{sec:ch3-ma-nguon-mo}

\textbf{Prometheus Alertmanager.} Theo tài liệu chính thức của dự án Prometheus~\cite{prometheus2026alertmanager}, Alertmanager đảm nhận tầng xử lý thông báo sau khi các luật cảnh báo Prometheus kích hoạt, cung cấp bốn cơ chế nền tảng: cây định tuyến theo nhãn (\code{route}), gom nhóm cảnh báo cùng nhãn thành một thông báo (\code{group\_by}), ức chế cảnh báo đích khi cảnh báo nguồn cùng nhóm đang kích hoạt (\code{inhibit\_rules}) và tắt tiếng tạm thời theo khung bảo trì (\code{silences}). Alertmanager nhẹ, minh bạch và là tiêu chuẩn thực tế trong hệ sinh thái đám mây, nhưng mọi quan hệ ức chế và gom nhóm đều phải cấu hình tĩnh bằng tay trong tệp YAML; công cụ không tự suy diễn đồ thị phụ thuộc từ trace, không đọc nội dung log và không thực hiện thẩm định cảnh báo hay phân tích nguyên nhân gốc.

\textbf{Keep.} Keep là nền tảng AIOps và quản lý cảnh báo mã nguồn mở; theo tài liệu chính thức của dự án~\cite{keep2026docs}, hệ thống hỗ trợ khử trùng lặp cảnh báo bằng dấu vân tay (\emph{fingerprint deduplication}, gồm chế độ toàn phần và từng phần), cung cấp ba cơ chế tương quan cảnh báo (luật thủ công bằng biểu thức CEL, tương quan theo topology và tương quan hỗ trợ bởi AI) cùng quy trình tự động hóa khai báo bằng YAML để làm giàu dữ liệu và gọi công cụ bên ngoài. Keep khắc phục được tính đơn điệu của công cụ luật tĩnh và mở rộng kết nối tốt với nhiều nguồn giám sát, nhưng vẫn xoay quanh bài toán tổng hợp--điều phối cảnh báo và quy trình tự động hóa; tài liệu hiện hành chưa đề cập cơ chế hiệu chỉnh \emph{LLM-as-a-Judge} theo nguồn bằng chứng hay khung thực nghiệm tái lập đánh giá đánh đổi recall--FPR dưới ngân sách LLM cố định.

\textbf{HolmesGPT.} HolmesGPT là tác tử điều tra sự cố mã nguồn mở khởi nguồn từ Robusta và hiện là dự án CNCF Sandbox; theo tài liệu chính thức~\cite{holmesgpt2026docs}, công cụ vận hành theo vòng lặp ReAct chạy ngay trong cụm Kubernetes ở chế độ chỉ đọc và tuân thủ phân quyền RBAC. Khi nhận cảnh báo từ Alertmanager hoặc PagerDuty, HolmesGPT để LLM tự chọn và gọi các công cụ truy vấn dữ liệu quan sát (như Prometheus, Loki hoặc \code{kubectl}), đọc kết quả trả về qua nhiều bước rồi tổng hợp nhận định nguyên nhân gốc. HolmesGPT minh họa rõ tính khả thi của tác tử LLM mã nguồn mở trong vận hành thực tế, nhưng chủ yếu đảm nhận khâu điều tra sau khi cảnh báo đã gửi tới, tốn nhiều lượt gọi LLM cho mỗi cảnh báo theo vòng lặp mở, và chưa tích hợp tầng thẩm định hiệu chỉnh trước đó để chặn báo động giả hay khung đánh giá thực nghiệm thống kê trên bộ chuẩn công khai.

\subsection{Nhận xét về các hệ thống hiện có}
\label{sec:ch3-nhan-xet}

Qua đối chiếu ba nền tảng thương mại và ba công cụ mã nguồn mở, có thể rút ra hai nhận xét chính về hiện trạng công nghiệp. Thứ nhất, các nền tảng thương mại (Datadog, Dynatrace, PagerDuty) đã tích hợp sâu khả năng gom cảnh báo theo topology và từng bước đưa tác tử AI vào điều tra sự cố, nhưng tất cả đều là hệ thống hộp đen gắn với hệ sinh thái thu thập độc quyền; chúng không cho phép kiểm định độc lập tác động của từng thành phần truy xuất hay chính sách hiệu chỉnh lên tỷ lệ báo động giả và không phục vụ được mục tiêu tái lập thực nghiệm khoa học. Thứ hai, ở mảng mã nguồn mở, hai năng lực cốt lõi vẫn nằm tách rời: nhóm quản lý cảnh báo (Alertmanager, Keep) mạnh về khử trùng lặp, luật ức chế và quy trình gom nhóm nhưng thiếu bộ thẩm định dựa trên bằng chứng log--trace đã hiệu chỉnh, còn tác tử điều tra bằng LLM (HolmesGPT) lại chạy trực tiếp theo vòng lặp nhiều bước mà không có tầng lọc kiểm soát ngân sách và kìm báo động giả ở đầu vào.

Do đó, thay vì cạnh tranh về độ phủ kết nối hạ tầng sản xuất với các sản phẩm thương mại, nền tảng trong đề tài này được xây dựng như một kiến trúc mở gọn nhẹ phục vụ đồng thời hai mục đích: tái lập chặt chẽ các thí nghiệm so sánh và loại trừ trên bộ dữ liệu công khai, và trình diễn trọn vẹn luồng xử lý từ phát lại dữ liệu quan sát, thẩm định cảnh báo bằng \CTAV{} đến định vị nguyên nhân gốc có bằng chứng.

% -----------------------------------------------------------------------------
\section{Kết chương}
\label{sec:ch3-ket-chuong}

Mục này tổng hợp các nghiên cứu học thuật ở Mục~\ref{sec:ch3-nghien-cuu} và các hệ thống thực tế ở Mục~\ref{sec:ch3-he-thong} trên cùng một hệ tiêu chí, qua đó chỉ rõ bốn khoảng trống nghiên cứu và định vị đóng góp của đề tài.

\subsection{Bảng so sánh tổng hợp}
\label{sec:ch3-bang-so-sanh}

Bảng~\ref{tab:ch3-so-sanh} đối chiếu các nhóm công trình tiêu biểu và các hệ thống hiện có theo năm tiêu chí bám sát trực tiếp mục tiêu của đề tài: (1) khai thác đồ thị topology dịch vụ, (2) sử dụng LLM để thẩm định hoặc suy luận sự cố, (3) hiệu chỉnh quyết định theo nguồn gốc bằng chứng, (4) thực hiện phân tích nguyên nhân gốc (RCA), và (5) có đánh giá định lượng báo động giả (FP/FPR) trên quy trình mở tái lập được. Trong bảng, \emph{Có} biểu thị tiêu chí được đáp ứng trực tiếp; \emph{Một phần} biểu thị tiêu chí chỉ được đáp ứng gián tiếp hoặc ở phạm vi hẹp (chẳng hạn chỉ có luật ức chế tĩnh, chỉ học đồ thị cảm biến không phải topology dịch vụ, hoặc hệ thống thương mại đóng nguồn không công bố kiểm định FP độc lập); \emph{Không} biểu thị tiêu chí nằm ngoài phạm vi của công trình hoặc công cụ.

\begin{table}[htbp]
  \centering
  \caption[So sánh công trình liên quan, hệ thống hiện có và đề xuất]{So sánh các nhóm nghiên cứu liên quan, hệ thống hiện có và đề xuất của đề tài.}
  \label{tab:ch3-so-sanh}
  \footnotesize
  \renewcommand{\arraystretch}{1.15}
  \setlength{\tabcolsep}{3.5pt}
  \begin{tabularx}{\linewidth}{@{}Y >{\centering\arraybackslash}p{1.45cm} >{\centering\arraybackslash}p{1.45cm} >{\centering\arraybackslash}p{1.6cm} >{\centering\arraybackslash}p{1.2cm} >{\centering\arraybackslash}p{1.65cm}@{}}
    \toprule
    \textbf{Công trình / Hệ thống} &
    \textbf{Topology} &
    \textbf{LLM judge} &
    \textbf{Hiệu chỉnh BC} &
    \textbf{RCA} &
    \textbf{Đánh giá FP \& tái lập} \\
    \midrule
    Isolation Forest~\cite{liu2008iforest}, RRCF~\cite{guha2016rrcf} &
      Không & Không & Không & Không & Một phần \\
    MTAD-GAT~\cite{zhao2020mtadgat}, GDN~\cite{deng2021gdn} &
      Một phần & Không & Không & Không & Một phần \\
    LogParser-LLM~\cite{zhong2024logparser}, LogLLM~\cite{guan2024logllm} &
      Không & Có & Không & Không & Một phần \\
    RCACopilot~\cite{chen2024rcacopilot}, RCAgent~\cite{wang2024rcagent}, Roy~\cite{roy2024exploring} &
      Không & Có & Không & Có & Không \\
    RAAD-LLM~\cite{russellgilbert2025raad}, ARGOS~\cite{gu2025argos}, CALM~\cite{devireddy2025calm} &
      Không & Có & Không & Không & Một phần \\
    MonitorRank~\cite{kim2013monitorrank}, MicroRCA~\cite{wu2020microrca}, CIRCA~\cite{li2022circa}, RCD~\cite{ikram2022rcd} &
      Có & Không & Không & Có & Không \\
    Nezha~\cite{yu2023nezha}, BARO~\cite{pham2024baro}, RCAEval~\cite{pham2025rcaeval} &
      Có & Không & Không & Có & Một phần \\
    GraphRAG~\cite{edge2024graphrag} &
      Một phần & Có & Không & Không & Không \\
    Truy xuất topology cho thẩm định LLM~\cite{author2026icaetm} &
      Có & Có & Không & Không & Có \\
    FairEval~\cite{wang2024fair}, Xiong~\cite{xiong2024uncertainty}, ALCE~\cite{gao2023alce} &
      Không & Có & Một phần & Không & Một phần \\
    Alert Storm~\cite{zhao2020alertstorm}, Kuang~\cite{kuang2024knowledge} &
      Một phần & Không & Không & Không & Một phần \\
    \midrule
    Datadog (Event Mgmt + Bits AI)~\cite{datadog2026docs} &
      Có & Có & Không & Có & Không \\
    Dynatrace (Davis AI + Smartscape)~\cite{dynatrace2026docs} &
      Có & Không & Không & Có & Không \\
    PagerDuty AIOps~\cite{pagerduty2026docs} &
      Một phần & Không & Không & Một phần & Không \\
    Prometheus Alertmanager~\cite{prometheus2026alertmanager} &
      Không & Không & Không & Không & Một phần \\
    Keep~\cite{keep2026docs} &
      Có & Một phần & Không & Không & Một phần \\
    HolmesGPT~\cite{holmesgpt2026docs} &
      Một phần & Có & Không & Có & Không \\
    \midrule
    \textbf{Đề tài này (\CTAV{} + tác tử RCA)} &
      \textbf{Có} & \textbf{Có} & \textbf{Có} & \textbf{Có} & \textbf{Có} \\
    \bottomrule
  \end{tabularx}
\end{table}

Bảng~\ref{tab:ch3-so-sanh} cho thấy bức tranh phân mảnh rõ rệt giữa các hướng hiện có. Nhóm phát hiện bất thường chuỗi thời gian và kiến trúc giám sát lai chưa khai thác cấu trúc đồ thị dịch vụ; nhóm RCA dựa trên đồ thị và nhân quả khai thác tốt topology nhưng không dùng LLM thẩm định luồng cảnh báo thô để kìm báo động giả; nhóm tác tử RCA dùng LLM bỏ qua khâu lọc cảnh báo giả có kiểm soát ngân sách; còn nhóm hiệu chỉnh \emph{LLM-as-a-Judge} mới xét bài toán ngôn ngữ tổng quát chứ chưa hiệu chỉnh theo nguồn gốc bằng chứng trên đồ thị phụ thuộc. Các hệ thống công nghiệp hoặc đóng nguồn không thể tái lập kiểm định FP, hoặc tách rời tầng quản lý cảnh báo với tầng điều tra LLM.

\subsection{Khoảng trống nghiên cứu và định vị đề xuất}
\label{sec:ch3-khoang-trong}

Từ kết quả khảo sát và Bảng~\ref{tab:ch3-so-sanh}, đề tài xác định bốn khoảng trống nghiên cứu cụ thể:

\begin{enumerate}
  \item \textbf{Khoảng trống G1 — Khoảng cách giữa phát hiện tốc độ cao và thẩm định chính xác dưới ngân sách LLM hữu hạn.} Bộ phát hiện metric không giám sát chạy nhanh nhưng dễ tạo báo động giả hoặc bỏ sót lỗi ngữ nghĩa; ngược lại, tác tử LLM gọi nhiều công cụ trực tiếp quá tốn kém để chạy trên mọi cửa sổ thời gian. Cần một cơ chế kết nối có kiểm soát chặt số lượt gọi LLM mà vẫn bảo đảm chất lượng thẩm định.
  \item \textbf{Khoảng trống G2 — Đánh đổi chưa được kiểm soát giữa thu hồi theo topology và báo động giả do thiên lệch xác nhận.} Mở rộng truy xuất bằng chứng sang láng giềng trên đồ thị phụ thuộc giúp nhận diện các lỗi lan truyền mà truy xuất từ khóa cục bộ bỏ lỡ, nhưng đồng thời khuếch đại xu hướng xác nhận nhầm của \emph{LLM-as-a-Judge} trên các cửa sổ bình thường nếu không có chính sách lọc và neo bằng chứng phù hợp.
  \item \textbf{Khoảng trống G3 — Thiếu cơ chế hiệu chỉnh bộ thẩm định cảnh báo theo nguồn gốc bằng chứng và thiếu kiểm chứng loại trừ có đăng ký giả thuyết trước.} Các phương pháp hiệu chỉnh LLM hiện hành chưa mô hình hóa sự khác biệt giữa bằng chứng nội bộ và bằng chứng láng giềng hay tham chiếu ngưỡng nền từ tập phát triển. Đồng thời, trong thực nghiệm AIOps, nhiều thiết kế phức tạp được gộp chung vào một hệ thống mà không tách kiểm định từng thành phần trên tập kiểm tra giữ kín, dễ che giấu những thành phần không mang lại lợi ích thực tế.
  \item \textbf{Khoảng trống G4 — Sự đứt gãy giữa thẩm định cảnh báo đã kiểm soát FP với phân tích nguyên nhân gốc có trích dẫn bằng chứng trên một nền tảng mở.} Phần lớn công trình hoặc dừng ở bước phát hiện/thẩm định cảnh báo, hoặc bắt đầu RCA từ danh sách sự cố đã xác nhận sẵn, và hiếm khi kết hợp cả hai trên một nền tảng mã nguồn mở có bộ nhớ đệm lời gọi LLM để tái lập thực nghiệm và trình diễn trực quan.
\end{enumerate}

Để lấp đầy bốn khoảng trống trên, đề tài đề xuất \textbf{\CTAV{}} (\emph{Calibrated Topology-Aware Alert Validation}) và tích hợp \CTAV{} làm tầng nền cho \textbf{tác tử phân tích nguyên nhân gốc} trên một nền tảng thống nhất. Cụ thể, trong không gian thiết kế của \CTAV{}, đề tài xây dựng và kiểm định tách biệt từng thành phần theo quy trình đăng ký giả thuyết trước: thành phần \textbf{C1} (truy xuất lai BM25--topology, đang dùng trong cấu hình cuối) kết hợp với thành phần \textbf{C3} (neo bằng chứng theo ngưỡng nền đo trên tập phát triển \splitname{dev}, đang dùng trong cấu hình cuối) đóng vai trò lõi giúp duy trì lợi ích nhận diện của topology đồng thời kìm mạnh báo động giả. Hai cơ chế bổ sung trong thiết kế ban đầu là \textbf{C2} (ngưỡng quyết định bất đối xứng theo nguồn bằng chứng, đã loại sau kiểm định) và \textbf{C4} (định tuyến hai tầng thích nghi, đã loại sau kiểm định) được đánh giá minh bạch trên tập \splitname{val} và loại khỏi cấu hình cuối vì không tạo lợi ích bổ sung so với cặp C1+C3 tinh gọn; tương tự, cơ chế \textbf{C5} (\code{Rescue Score}, chỉ đánh giá ở vai trò phân tích loại trừ) được khảo sát độc lập thay vì mặc định bật trong hệ thống chính. Phía sau \CTAV{}, tác tử RCA gom các cảnh báo đã thẩm định thành sự cố, lần ngược đường lan truyền trên đồ thị phụ thuộc dịch vụ và sinh báo cáo chẩn đoán có cấu trúc kèm trích dẫn bằng chứng. Toàn bộ thiết kế kiến trúc, thuật toán chi tiết và mô hình dữ liệu của hệ thống được trình bày trong Chương~\ref{ch:de-xuat}.

%        ...
%        \bibliographystyle{IEEEtran}
%        \bibliography{refs}
%
% Cac nhan lien chuong duoc tham chieu tu Chuong 3:
%   Chuong: ch:mo-dau, ch:nen-tang, ch:de-xuat, ch:thuc-nghiem, ch:platform
%   Chuong 2: sec:ch2-du-lieu, sec:ch2-truy-xuat, sec:ch2-judge, sec:ch2-rca,
%             sec:ch2-dang-ky
%   Chuong 4: sec:ch4-ctav, sec:ch4-c1, sec:ch4-policy, sec:ch4-c4, sec:ch4-rca
%   Chuong 5: sec:ch5-hypotheses, sec:ch5-c5, sec:ch5-budget
% =============================================================================

\chapter{Công trình liên quan}
\label{ch:cong-trinh-lien-quan}

Chương này hệ thống hóa các hướng nghiên cứu và các hệ thống vận hành thực tế liên quan trực tiếp tới bài toán thẩm định cảnh báo có nhận biết topology và phân tích nguyên nhân gốc trên kiến trúc microservice. Nội dung được tổ chức xoay quanh hai câu hỏi: cộng đồng học thuật và công nghiệp đã giải quyết những khía cạnh nào của chuỗi xử lý từ phát hiện bất thường đến định vị nguyên nhân gốc, và còn những khoảng trống nào mà đề xuất của đề tài nhắm tới. Mục~\ref{sec:ch3-nghien-cuu} khảo sát bảy nhánh nghiên cứu học thuật từ phát hiện bất thường chuỗi thời gian, ứng dụng mô hình ngôn ngữ lớn trong phân tích log và nguyên nhân gốc, kiến trúc giám sát lai, phân tích nguyên nhân gốc dựa trên đồ thị, truy xuất dựa trên đồ thị, hiệu chỉnh độ tin cậy của \emph{LLM-as-a-Judge}, đến gom nhóm và phân loại cảnh báo. Mục~\ref{sec:ch3-he-thong} đối chiếu các nền tảng AIOps thương mại và công cụ mã nguồn mở hiện hành dựa trên tài liệu chính thức. Cuối cùng, Mục~\ref{sec:ch3-ket-chuong} tổng hợp các hướng tiếp cận trong Bảng~\ref{tab:ch3-so-sanh}, rút ra bốn khoảng trống nghiên cứu và định vị rõ phạm vi đóng góp của \CTAV{} cùng tác tử phân tích nguyên nhân gốc.

% -----------------------------------------------------------------------------
\section{Các nghiên cứu liên quan}
\label{sec:ch3-nghien-cuu}

Các công trình liên quan đến đề tài nằm ở giao điểm giữa giám sát chuỗi thời gian, phân tích dữ liệu quan sát đa phương thức, truy xuất bằng chứng theo đồ thị và hiệu chỉnh mô hình ngôn ngữ lớn đóng vai trò bộ thẩm định. Bảy tiểu mục dưới đây phân tích từng nhánh theo cùng một trục: cách tiếp cận của công trình, hạn chế đặt trong bài toán của đề tài và điểm khác biệt của đề tài.

\subsection{Phát hiện bất thường trên chuỗi thời gian metric}
\label{sec:ch3-tsad}

Hướng phát hiện bất thường trên chuỗi thời gian không giám sát hình thành nền tảng cho tầng lọc cảnh báo đầu vào. Liu và cộng sự~\cite{liu2008iforest} đề xuất Isolation Forest, một phương pháp dựa trên tập cây phân tách ngẫu nhiên tận dụng nhận xét rằng điểm bất thường ít và khác biệt nên cần ít phép chia để cô lập hơn điểm bình thường; mô hình chạy nhanh, không cần nhãn và tiêu thụ ít tài nguyên. Mở rộng sang dữ liệu luồng có phân bố biến đổi theo thời gian, Guha và cộng sự~\cite{guha2016rrcf} xây dựng Robust Random Cut Forest (RRCF) cho phép cập nhật cấu trúc cây trực tuyến khi điểm mới đến và đo độ bất thường bằng mức thay đổi độ phức tạp cây khi chèn hoặc xóa điểm. Hạn chế chung của nhóm cây phân tách là chỉ nhìn từng vector metric độc lập theo cửa sổ, nên không phân biệt được dao động tải bình thường với sự cố thực sự nếu không đối chiếu ngữ cảnh văn bản từ log hoặc quan hệ gọi giữa các service.

Để nắm bắt tương quan đa biến và phụ thuộc giữa các cảm biến hay metric, các công trình học sâu dựa trên đồ thị đã được đề xuất. Zhao và cộng sự~\cite{zhao2020mtadgat} giới thiệu MTAD-GAT với hai lớp chú ý đồ thị song song: một lớp học tương quan giữa các đặc trưng tại cùng thời điểm và một lớp học phụ thuộc dọc trục thời gian, kết hợp đồng thời mục tiêu dự báo và tái tạo để tính điểm bất thường. Deng và Hooi~\cite{deng2021gdn} đề xuất Graph Deviation Network (GDN), tự học đồ thị quan hệ giữa các cảm biến qua vector nhúng và dùng mạng chú ý đồ thị để dự báo giá trị tiếp theo, từ đó giải thích bất thường qua độ lệch của các cạnh dự báo. Tuy đạt độ chính xác cao hơn trên nhiều tập chuỗi thời gian đa biến, nhóm học sâu đồ thị đòi hỏi huấn luyện lại khi topology hoặc tập metric thay đổi, cần tài nguyên tính toán lớn hơn đáng kể và vẫn chỉ khai thác tín hiệu số mà bỏ qua nội dung log hay trace vốn chứa thông tin tường minh về lỗi ứng dụng.

Khác với các công trình tập trung thay thế bộ phát hiện metric bằng kiến trúc học sâu ngày càng phức tạp, đề tài này giữ tầng phát hiện nhẹ bằng Isolation Forest để bảo đảm tốc độ và khả năng chạy lại tức thì trên mọi cấu hình, rồi chuyển trọng tâm sang tầng thẩm định bằng bằng chứng đa phương thức phía sau. Cách phân công này cho phép tận dụng tốc độ lọc sơ bộ của detector nhẹ trong khi vẫn kiểm soát được báo động giả nhờ đối chiếu log và đồ thị phụ thuộc service.

\subsection{Mô hình ngôn ngữ lớn trong phân tích log và nguyên nhân gốc}
\label{sec:ch3-llm-log-rca}

Sự xuất hiện của các mô hình ngôn ngữ lớn (LLM) mở ra hướng xử lý trực tiếp dữ liệu log phi cấu trúc và suy luận sự cố bằng ngôn ngữ tự nhiên. Ở bài toán phân tách mẫu log (\emph{log parsing}), Zhong và cộng sự~\cite{zhong2024logparser} đề xuất LogParser-LLM, kết hợp cây tiền tố với truy vấn LLM chọn lọc để trích xuất mẫu log mà không cần huấn luyện lại mô hình khi định dạng log thay đổi. Ở bài toán phát hiện bất thường trên chuỗi log, Guan và cộng sự~\cite{guan2024logllm} giới thiệu LogLLM, phối hợp bộ mã hóa ngữ nghĩa dạng BERT cho từng dòng log với bộ giải mã LLM nhằm mô hình hóa chuỗi log mà không cần bước tách mẫu thủ công. Hai công trình này giải quyết tốt khâu biểu diễn và phân tích log đơn lẻ, nhưng chỉ làm việc trên tín hiệu log tách rời mà không tích hợp luồng cảnh báo từ metric cũng như cấu trúc phụ thuộc giữa các service.

Ở bài toán phân tích nguyên nhân gốc cho sự cố đám mây, nhiều công trình sử dụng LLM trực tiếp trên hồ sơ sự cố hoặc dưới dạng tác tử gọi công cụ. Ahmed và cộng sự~\cite{ahmed2023recommending} thực hiện đánh giá quy mô lớn trên hơn 40.000 sự cố vận hành thực tế tại Microsoft để khuyến nghị nguyên nhân gốc và bước khắc phục bằng mô hình dòng GPT, cho thấy LLM có ích trong việc tóm tắt tri thức sự cố nhưng nhạy cảm với chất lượng mô tả đầu vào. Chen và cộng sự~\cite{chen2024rcacopilot} xây dựng RCACopilot, một hệ thống định tuyến sự cố vào quy trình thu thập dữ liệu chẩn đoán tương ứng rồi để LLM dự đoán nhóm nguyên nhân gốc kèm giải thích bằng ngôn ngữ tự nhiên. Tiến thêm bước tự chủ, Wang và cộng sự~\cite{wang2024rcagent} đề xuất RCAgent sử dụng tác tử LLM có hỗ trợ công cụ truy vấn log và tổng hợp tự nhất quán (\emph{self-consistency}) trong môi trường đám mây, trong khi Roy và cộng sự~\cite{roy2024exploring} khảo sát tác tử theo vòng lặp ReAct trang bị công cụ truy xuất log, metric và cơ sở dữ liệu lịch sử để điều tra sự cố.

Điểm hạn chế chung của nhóm công trình LLM cho RCA nêu trên là phần lớn giả định sự cố đầu vào đã được xác nhận hoặc đã mở phiếu sự cố sẵn, nên không đánh giá trực tiếp bài toán thẩm định luồng cảnh báo thô để chặn báo động giả trước khi gọi tác tử chẩn đoán tốn kém. Ngoài ra, các hệ thống đánh giá trên dữ liệu sự cố nội bộ doanh nghiệp khó tái lập độc lập bên ngoài. Đề tài này đặt tầng thẩm định cảnh báo có kiểm soát ngân sách ngay trước tác tử RCA và thực nghiệm hoàn toàn trên bộ dữ liệu công khai có thể tái lập.

\subsection{Kiến trúc giám sát lai và LLM-as-a-Judge}
\label{sec:ch3-hybrid-judge}

Để dung hòa chi phí gọi LLM với yêu cầu giám sát liên tục, nhiều công trình kết hợp mô hình nhẹ hoặc cơ chế truy xuất tăng cường sinh (\emph{Retrieval-Augmented Generation}, RAG)~\cite{lewis2020rag} với bộ phân xử LLM ở tầng sau. Về nguyên lý phân tầng chi phí, nền tảng lý thuyết bắt nguồn từ cơ chế đánh đổi sai số và quyền từ chối (\emph{reject option}) của Chow~\cite{chow1970reject} cùng mô hình chuỗi tầng gọi LLM tiết kiệm chi phí FrugalGPT của Chen và cộng sự~\cite{chen2024frugalgpt}, trong đó các mẫu rõ ràng được xử lý bởi mô hình rẻ hơn và chỉ vùng bất định mới chuyển lên mô hình đắt hơn.

Trong miền phát hiện bất thường, một loạt công trình gần đây đưa kiến trúc lai vào thực nghiệm. Russell-Gilbert và cộng sự~\cite{russellgilbert2025raad} đề xuất RAAD-LLM, tích hợp RAG vào LLM để thích nghi phát hiện bất thường chuỗi thời gian khi bối cảnh hệ thống thay đổi mà không cần tinh chỉnh trọng số. Gu và cộng sự~\cite{gu2025argos} giới thiệu ARGOS, một khuôn khổ tác tử cho phép LLM tự sinh và tinh chỉnh luật phát hiện bất thường trên chuỗi thời gian nhằm giảm phụ thuộc vào truy vấn trực tiếp trên từng mẫu. Devireddy và Huang~\cite{devireddy2025calm} đề xuất CALM, điều phối liên tục giữa phát hiện thời gian thực và cơ chế can thiệp thích nghi của LLM. Ở góc độ đánh giá năng lực thẩm định, Sivalingam và cộng sự~\cite{sivalingam2026tsquery} xây dựng TSQueryBench để kiểm tra khả năng dùng \emph{LLM-as-a-Judge} đánh giá lời giải thích bất thường chuỗi thời gian.

Hạn chế chung của các kiến trúc lai này khi áp vào microservice là chúng chủ yếu vận hành trên chuỗi thời gian đơn lẻ hoặc truy xuất ngữ cảnh văn bản phẳng, chưa kết nối tường minh với đồ thị phụ thuộc giữa các service sinh từ trace. Quan trọng hơn, các nghiên cứu này chưa khảo sát hệ quả ngược chiều khi mở rộng bằng chứng cho bộ thẩm định LLM: việc cấp thêm ngữ cảnh từ láng giềng có thể làm tăng tỷ lệ báo động giả do thiên lệch xác nhận. Đề tài này tập trung giải quyết trực tiếp bài toán kiểm soát ngân sách và kìm báo động giả trong kiến trúc lai detector nhẹ kết hợp \emph{LLM-as-a-Judge} cho hệ microservice.

\subsection{Phân tích nguyên nhân gốc dựa trên topology, đồ thị và suy luận nhân quả}
\label{sec:ch3-rca-graph}

Định vị nguyên nhân gốc dựa trên đồ thị phụ thuộc và suy luận nhân quả là nhánh phát triển lâu đời nhất trong giám sát hệ phân tán. Kim và cộng sự~\cite{kim2013monitorrank} đề xuất MonitorRank, thực hiện bước đi ngẫu nhiên (\emph{random walk}) có trọng số tương quan metric trên đồ thị gọi dịch vụ để xếp hạng service nghi vấn trong kiến trúc hướng dịch vụ. Chen và cộng sự~\cite{chen2014causeinfer} xây dựng CauseInfer với đồ thị nhân quả hai tầng (tầng dịch vụ và tầng metric) bằng thuật toán kiểm định độc lập có điều kiện PC để lần ngược nút gốc gây suy giảm hiệu năng. Trên kiến trúc microservice, Wu và cộng sự~\cite{wu2020microrca} giới thiệu MicroRCA, gắn dữ liệu sử dụng tài nguyên máy chủ vào đồ thị dịch vụ bất thường và áp dụng Personalized PageRank để định vị service gây lỗi hiệu năng.

Các công trình gần đây tiếp cận bài toán bằng mô hình nhân quả chặt chẽ hơn và dữ liệu đa phương thức. Li và cộng sự~\cite{li2022circa} đề xuất CIRCA, mô hình hóa lỗi như một can thiệp (\emph{intervention}) trên mạng Bayes nhân quả xây từ kiến trúc hệ thống và dùng điều kiện thay đổi phân bố có điều kiện để chỉ ra chỉ số gốc. Ikram và cộng sự~\cite{ikram2022rcd} đề xuất RCD, xem sự cố như can thiệp cục bộ và khám phá đồ thị nhân quả theo phân cấp để tránh chi phí tính toán lớn của thuật toán PC trên toàn bộ hệ thống. Kết hợp cả ba trụ cột dữ liệu quan sát, Yu và cộng sự~\cite{yu2023nezha} đề xuất Nezha, chuyển đổi metric, log và trace thành biểu diễn sự kiện đồng nhất và so sánh mẫu phân kỳ giữa trạng thái bình thường với trạng thái lỗi để định vị nguyên nhân gốc ở mức chi tiết, đồng thời công bố bộ dữ liệu microservice đa phương thức làm chuẩn đánh giá. Đáng lưu ý về tính bền vững của chuỗi xử lý, Pham và cộng sự~\cite{pham2024baro} đề xuất BARO bằng phát hiện điểm đổi Bayes đa biến kết hợp kiểm định phi tham số, đồng thời chỉ ra rằng sai số ở khâu phát hiện bất thường đầu vào có thể lan sang làm giảm mạnh độ chính xác định vị gốc; tiếp đó nhóm tác giả công bố bộ chuẩn RCAEval~\cite{pham2025rcaeval} tổng hợp 735 ca lỗi đa phương thức trên nhiều hệ microservice mã nguồn mở để thống nhất việc so sánh.

Hạn chế của nhóm phương pháp đồ thị và nhân quả thuần túy gồm hai điểm. Thứ nhất, phần lớn mô hình chỉ xuất ra danh sách xếp hạng service hoặc metric gốc dưới dạng chỉ số toán học, thiếu báo cáo tường minh giải thích cơ chế lan truyền bằng chứng từ log và trace cho kỹ sư vận hành. Thứ hai, trừ những phân tích gần đây như BARO, đa số nghiên cứu RCA giả định cửa sổ thời gian sự cố đã được xác định sạch từ trước, tách rời khỏi bài toán lọc cảnh báo giả trực tuyến. Đề tài này nối trực tiếp hai khâu: dùng \CTAV{} để làm sạch luồng cảnh báo trước khi gom sự cố, rồi dùng tác tử RCA lần ngược chuỗi lan truyền trên đồ thị phụ thuộc và sinh báo cáo có cấu trúc kèm trích dẫn bằng chứng.

\subsection{Truy xuất dựa trên đồ thị và GraphRAG cho vận hành hệ thống}
\label{sec:ch3-graphrag}

Trong lĩnh vực xử lý ngôn ngữ tự nhiên, Edge và cộng sự~\cite{edge2024graphrag} đề xuất GraphRAG nhằm khắc phục hạn chế của truy xuất văn bản phẳng khi cần tổng hợp thông tin phân tán trên nhiều tài liệu. Phương pháp xây dựng đồ thị tri thức thực thể--quan hệ từ kho văn bản, phân cụm cộng đồng đồ thị và sinh tóm tắt phân cấp để trả lời các truy vấn mang tính toàn cục. Tuy mở ra hướng truy xuất có cấu trúc, GraphRAG thiết kế cho kho tài liệu tĩnh với đồ thị trích xuất bằng LLM, nên vừa tốn chi phí lập chỉ mục vừa khác về bản chất so với đồ thị phụ thuộc động của hệ microservice được đo trực tiếp từ quan hệ gọi \code{SpanID}--\code{ParentID} trong trace.

Khi áp dụng tư tưởng truy xuất theo cấu trúc đồ thị vào bài toán thẩm định cảnh báo microservice, một nghiên cứu tiền đề về truy xuất bằng chứng có nhận biết topology~\cite{author2026icaetm} đã đề xuất kết hợp điểm từ khóa BM25 với độ gần về thời gian và khoảng cách bước nhảy trên đồ thị phụ thuộc service để cung cấp thêm bằng chứng từ các service láng giềng cho bộ thẩm định LLM (chi tiết công thức hợp nhất đã trình bày ở Mục~\ref{sec:ch2-truy-xuat}). Kết quả định tính của hướng tiếp cận này cho thấy một đánh đổi rõ rệt: việc mở rộng ngữ cảnh sang các service có liên kết gọi giúp cải thiện khả năng nhận diện các sự cố lan truyền mà bằng chứng cục bộ không đủ, nhưng khi đưa trực tiếp vào \emph{LLM-as-a-Judge} mà không có cơ chế kiểm soát chính sách, số báo động giả tăng lên mạnh do bộ thẩm định dễ diễn giải nhầm log nền từ láng giềng thành dấu hiệu bất thường. Đề tài này xuất phát trực tiếp từ đánh đổi đó để xây dựng cơ chế hiệu chỉnh theo nguồn gốc bằng chứng và kiểm định nghiêm ngặt hiệu quả kìm báo động giả trên dữ liệu ngoài tập tinh chỉnh.

\subsection{Hiệu chỉnh độ tin cậy, thiên lệch của LLM-as-a-Judge và gắn bằng chứng}
\label{sec:ch3-calibration}

Dù \emph{LLM-as-a-Judge} đạt mức đồng thuận cao với người đánh giá trên nhiều tác vụ tổng quát~\cite{zheng2023judging}, hàng loạt nghiên cứu đã chứng minh mô hình ngôn ngữ mắc nhiều thiên lệch có hệ thống và thường tự tin thái quá. Wang và cộng sự~\cite{wang2024fair} chỉ ra rằng bộ thẩm định LLM chịu thiên lệch vị trí nghiêm trọng khi so sánh cặp câu trả lời, dẫn đến thứ hạng dễ bị thao túng chỉ bằng cách đảo thứ tự trình bày, và đề xuất khung FairEval hiệu chỉnh bằng cách buộc sinh bằng chứng trước khi chấm điểm kết hợp trung bình hóa vị trí. Sharma và cộng sự~\cite{sharma2024sycophancy} phân tích hiện tượng chiều lòng (\emph{sycophancy}), cho thấy mô hình đã qua tinh chỉnh theo phản hồi con người có xu hướng đồng thuận với giả định ngầm trong câu hỏi ngay cả khi giả định đó sai — một cơ chế tương đồng trực tiếp với thiên lệch xác nhận khi bộ thẩm định nhận thông báo rằng ``cảnh báo đang kích hoạt''.

Về hiệu chỉnh xác suất và ước lượng độ bất định, Guo và cộng sự~\cite{guo2017calibration} chứng minh các mạng nơ-ron hiện đại tuy chính xác nhưng hiệu chỉnh kém và đề xuất phương pháp hiệu chỉnh hậu xử lý bằng tỉ lệ nhiệt độ (\emph{temperature scaling}). Trên LLM, Kadavath và cộng sự~\cite{kadavath2022mostly} khảo sát khả năng mô hình tự đánh giá xác suất đúng của phát biểu do chính nó sinh ra, cho thấy tín hiệu nội tại có thể hiệu chỉnh được trong câu hỏi trắc nghiệm nhưng khó khái quát khi chuyển phân bố tác vụ. Gần hơn với bài toán sử dụng API thương mại không truy cập được xác suất token thô, Xiong và cộng sự~\cite{xiong2024uncertainty} đánh giá thực nghiệm độ tin cậy bằng lời (\emph{verbalized confidence}) trên nhiều dòng LLM và chỉ ra mô hình thường gán điểm tin cậy quá cao, tập trung dày đặc ở vùng sát $1{,}0$. Song song đó, để hạn chế suy diễn vô căn cứ trong văn bản sinh ra, Gao và cộng sự~\cite{gao2023alce} đề xuất chuẩn đánh giá ALCE đo đồng thời tính đúng đắn và chất lượng trích dẫn nguồn bằng chứng trong câu trả lời của LLM.

Khoảng trống còn lại trong nhánh nghiên cứu này là các kỹ thuật hiệu chỉnh hiện có chủ yếu thiết kế cho bài toán chấm điểm cặp câu trả lời, hỏi--đáp mở hoặc sinh văn bản có trích dẫn tổng quát. Chưa có cơ chế nào hiệu chỉnh trực tiếp quyết định của \emph{LLM-as-a-Judge} trong bài toán thẩm định cảnh báo microservice theo \emph{nguồn gốc bằng chứng} (nội bộ hay láng giềng trên đồ thị phụ thuộc) và theo \emph{mức độ bất thường nền đo trên tập phát triển}. Đề tài này cụ thể hóa ý tưởng hiệu chỉnh bằng chính sách neo bằng chứng và kiểm định tách biệt từng thành phần thiết kế (ngưỡng bất đối xứng theo nguồn bằng chứng, neo ngưỡng nền trên tập phát triển và định tuyến hai tầng) theo quy trình đăng ký giả thuyết trước ở Mục~\ref{sec:ch2-dang-ky}.

\subsection{Gom nhóm, phân loại và cứu cảnh báo trong vận hành}
\label{sec:ch3-alert-triage}

Trong vận hành hệ thống trực tuyến quy mô lớn, xử lý bão cảnh báo và định tuyến sự cố là bài toán sống còn để giảm tải cho kỹ sư trực. Zhao và cộng sự~\cite{zhao2020alertstorm} khảo sát hiện tượng bão cảnh báo trong hệ dịch vụ trực tuyến công nghiệp và đề xuất quy trình hai bước gồm phát hiện thời điểm bão và gom nhóm các cảnh báo tương quan về thời gian--ngữ nghĩa thành một cụm sự cố đại diện. Ở khâu định tuyến sự cố, Chen và cộng sự~\cite{chen2019empirical,chen2019continuous} phân tích nguyên nhân định tuyến nhầm sự cố tại các dịch vụ đám mây lớn và xây dựng phương pháp phân loại liên tục tận dụng thông tin cập nhật trong quá trình thảo luận của kỹ sư. Gần đây hơn, Kuang và cộng sự~\cite{kuang2024knowledge} đề xuất phương pháp gom cảnh báo có nhận biết tri thức cho hệ đám mây quy mô lớn, kết hợp đặc trưng văn bản với tri thức vận hành và topology để gộp cảnh báo cùng gốc tốt hơn các luật tĩnh.

Các công trình gom nhóm và phân loại này giải quyết tốt khâu \emph{sau khi cảnh báo đã phát sinh}, giúp nén số lượng thông báo gửi tới người trực, nhưng thường không thẩm định lại bản thân cảnh báo đầu vào có phải là báo động giả do bộ phát hiện metric gây ra hay không. Mặt khác, một câu hỏi ít được kiểm chứng nghiêm ngặt là liệu các cửa sổ bị bộ phát hiện chính bỏ sót có thể được ``cứu'' lại bằng điểm phụ trợ từ tín hiệu log/trace (cơ chế \code{Rescue Score}) mà không làm vỡ tỷ lệ báo động giả hay không. Đề tài này tách rõ hai vấn đề: sử dụng gom nhóm theo cửa sổ và topology ở đầu vào tác tử RCA để hợp nhất các cảnh báo đã được \CTAV{} thẩm định thành sự cố, đồng thời đánh giá cơ chế cứu cảnh báo như một phân tích loại trừ độc lập (Mục~\ref{sec:ch5-c5}) thay vì mặc định đưa vào hệ thống chính.

% -----------------------------------------------------------------------------
\section{Các hệ thống và công cụ hiện có}
\label{sec:ch3-he-thong}

Bên cạnh các công trình học thuật, nhiều nền tảng AIOps thương mại và công cụ mã nguồn mở đã được triển khai rộng rãi trong thực tế vận hành. Mục này mô tả tính năng cốt lõi của các hệ thống tiêu biểu dựa trên tài liệu chính thức của nhà phát triển, sau đó chỉ ra điểm khác biệt về mục tiêu và giới hạn khi đặt trong bối cảnh nghiên cứu của đề tài.

\subsection{Nền tảng AIOps thương mại}
\label{sec:ch3-thuong-mai}

\textbf{Datadog (Event Management và Bits AI).} Theo tài liệu chính thức của Datadog~\cite{datadog2026docs}, phân hệ Event Management thu nạp tập trung cảnh báo từ các bộ giám sát nội bộ cùng nguồn bên thứ ba, sau đó áp dụng tương quan cảnh báo tự động dựa trên topology dịch vụ, nhãn chia sẻ và mẫu thời gian để gom các cảnh báo lan truyền vào chung một sự vụ, đồng thời làm giàu ngữ cảnh về đơn vị sở hữu dịch vụ. Ở lớp điều tra sự cố, tác tử Bits AI (Bits Investigation) được kích hoạt tự động khi có cảnh báo nghiêm trọng, thực hiện kiểm định song song nhiều giả thuyết trên dữ liệu metric, log, trace và thay đổi cấu hình để lần theo chuỗi nhân quả và đề xuất hướng xử lý. Tuy cung cấp quy trình khép kín từ gom cảnh báo đến điều tra bằng AI, kiến trúc bên trong của cơ chế tương quan và điều tra đều đóng kín; nhà cung cấp không công bố phương pháp hiệu chỉnh độ tin cậy của bộ suy luận LLM theo nguồn bằng chứng hay đánh giá tách biệt tác động của topology lên tỷ lệ báo động giả trên bộ chuẩn công khai.

\textbf{Dynatrace (Davis AI và Smartscape).} Tài liệu kỹ thuật của Dynatrace~\cite{dynatrace2026docs} mô tả kiến trúc dựa trên hai trụ cột là bản đồ phụ thuộc thời gian thực Smartscape và động cơ phân tích nhân quả Davis AI. Smartscape tự động lập đồ thị phụ thuộc đầy đủ từ ứng dụng, service, tiến trình đến máy chủ và hạ tầng Kubernetes. Khi phát sinh bất thường ở mức chỉ số, độ trễ hoặc tỷ lệ lỗi, Davis AI gom các sự kiện liên quan có cùng nguyên nhân gốc thành một thực thể \emph{Problem} duy nhất và thực hiện phân tích cây lỗi tất định (\emph{deterministic fault-tree analysis}) trên đồ thị Smartscape để tách biệt thành phần gốc khỏi các thành phần chịu ảnh hưởng, hiển thị qua biểu đồ \emph{Visual Resolution Path}. Ưu điểm của Dynatrace là tính nhất quán cao nhờ mô hình nhân quả dựa trên agent thu thập sâu (\emph{OneAgent}), nhưng hệ thống gắn chặt với hạ tầng thu thập độc quyền, chi phí bản quyền lớn và không nhằm giải bài toán hiệu chỉnh quyết định của \emph{LLM-as-a-Judge} trên dữ liệu log phi cấu trúc lẫn ngân sách truy vấn LLM giới hạn.

\textbf{PagerDuty AIOps.} Theo tài liệu hỗ trợ chính thức của PagerDuty~\cite{pagerduty2026docs}, bộ công cụ PagerDuty AIOps tập trung vào ba chức năng: \emph{Intelligent Alert Grouping} dùng học máy nhóm các cảnh báo có liên quan theo mẫu lịch sử để giảm nhiễu thông báo; \emph{Probable Origin} phân tích chuỗi sự cố lan truyền trong lịch sử sáu tháng gần nhất để gợi ý tối đa ba dịch vụ hoặc sự cố có khả năng là điểm khởi phát cao nhất; và \emph{Event Orchestration} cung cấp bộ máy luật cho phép lọc, làm giàu và định tuyến sự kiện tự động ở mức toàn cục hoặc từng dịch vụ. Cách tiếp cận của PagerDuty giải quyết hiệu quả bài toán quản trị phản ứng sự cố ở tầng thông báo, nhưng phụ thuộc nhiều vào lịch sử sự cố tích lũy theo tổ chức, không khai thác đồ thị gọi chi tiết từng cửa sổ trace để thẩm định cảnh báo từ tín hiệu thô và không công bố cơ chế kiểm soát báo động giả cho tác tử thẩm định LLM.

\subsection{Công cụ mã nguồn mở}
\label{sec:ch3-ma-nguon-mo}

\textbf{Prometheus Alertmanager.} Theo tài liệu chính thức của dự án Prometheus~\cite{prometheus2026alertmanager}, Alertmanager đảm nhận tầng xử lý thông báo sau khi các luật cảnh báo Prometheus kích hoạt, cung cấp bốn cơ chế nền tảng: cây định tuyến theo nhãn (\code{route}), gom nhóm cảnh báo cùng nhãn thành một thông báo (\code{group\_by}), ức chế cảnh báo đích khi cảnh báo nguồn cùng nhóm đang kích hoạt (\code{inhibit\_rules}) và tắt tiếng tạm thời theo khung bảo trì (\code{silences}). Alertmanager nhẹ, minh bạch và là tiêu chuẩn thực tế trong hệ sinh thái đám mây, nhưng mọi quan hệ ức chế và gom nhóm đều phải cấu hình tĩnh bằng tay trong tệp YAML; công cụ không tự suy diễn đồ thị phụ thuộc từ trace, không đọc nội dung log và không thực hiện thẩm định cảnh báo hay phân tích nguyên nhân gốc.

\textbf{Keep.} Keep là nền tảng AIOps và quản lý cảnh báo mã nguồn mở; theo tài liệu chính thức của dự án~\cite{keep2026docs}, hệ thống hỗ trợ khử trùng lặp cảnh báo bằng dấu vân tay (\emph{fingerprint deduplication}, gồm chế độ toàn phần và từng phần), cung cấp ba cơ chế tương quan cảnh báo (luật thủ công bằng biểu thức CEL, tương quan theo topology và tương quan hỗ trợ bởi AI) cùng quy trình tự động hóa khai báo bằng YAML để làm giàu dữ liệu và gọi công cụ bên ngoài. Keep khắc phục được tính đơn điệu của công cụ luật tĩnh và mở rộng kết nối tốt với nhiều nguồn giám sát, nhưng vẫn xoay quanh bài toán tổng hợp--điều phối cảnh báo và quy trình tự động hóa; tài liệu hiện hành chưa đề cập cơ chế hiệu chỉnh \emph{LLM-as-a-Judge} theo nguồn bằng chứng hay khung thực nghiệm tái lập đánh giá đánh đổi recall--FPR dưới ngân sách LLM cố định.

\textbf{HolmesGPT.} HolmesGPT là tác tử điều tra sự cố mã nguồn mở khởi nguồn từ Robusta và hiện là dự án CNCF Sandbox; theo tài liệu chính thức~\cite{holmesgpt2026docs}, công cụ vận hành theo vòng lặp ReAct chạy ngay trong cụm Kubernetes ở chế độ chỉ đọc và tuân thủ phân quyền RBAC. Khi nhận cảnh báo từ Alertmanager hoặc PagerDuty, HolmesGPT để LLM tự chọn và gọi các công cụ truy vấn dữ liệu quan sát (như Prometheus, Loki hoặc \code{kubectl}), đọc kết quả trả về qua nhiều bước rồi tổng hợp nhận định nguyên nhân gốc. HolmesGPT minh họa rõ tính khả thi của tác tử LLM mã nguồn mở trong vận hành thực tế, nhưng chủ yếu đảm nhận khâu điều tra sau khi cảnh báo đã gửi tới, tốn nhiều lượt gọi LLM cho mỗi cảnh báo theo vòng lặp mở, và chưa tích hợp tầng thẩm định hiệu chỉnh trước đó để chặn báo động giả hay khung đánh giá thực nghiệm thống kê trên bộ chuẩn công khai.

\subsection{Nhận xét về các hệ thống hiện có}
\label{sec:ch3-nhan-xet}

Qua đối chiếu ba nền tảng thương mại và ba công cụ mã nguồn mở, có thể rút ra hai nhận xét chính về hiện trạng công nghiệp. Thứ nhất, các nền tảng thương mại (Datadog, Dynatrace, PagerDuty) đã tích hợp sâu khả năng gom cảnh báo theo topology và từng bước đưa tác tử AI vào điều tra sự cố, nhưng tất cả đều là hệ thống hộp đen gắn với hệ sinh thái thu thập độc quyền; chúng không cho phép kiểm định độc lập tác động của từng thành phần truy xuất hay chính sách hiệu chỉnh lên tỷ lệ báo động giả và không phục vụ được mục tiêu tái lập thực nghiệm khoa học. Thứ hai, ở mảng mã nguồn mở, hai năng lực cốt lõi vẫn nằm tách rời: nhóm quản lý cảnh báo (Alertmanager, Keep) mạnh về khử trùng lặp, luật ức chế và quy trình gom nhóm nhưng thiếu bộ thẩm định dựa trên bằng chứng log--trace đã hiệu chỉnh, còn tác tử điều tra bằng LLM (HolmesGPT) lại chạy trực tiếp theo vòng lặp nhiều bước mà không có tầng lọc kiểm soát ngân sách và kìm báo động giả ở đầu vào.

Do đó, thay vì cạnh tranh về độ phủ kết nối hạ tầng sản xuất với các sản phẩm thương mại, nền tảng trong đề tài này được xây dựng như một kiến trúc mở gọn nhẹ phục vụ đồng thời hai mục đích: tái lập chặt chẽ các thí nghiệm so sánh và loại trừ trên bộ dữ liệu công khai, và trình diễn trọn vẹn luồng xử lý từ phát lại dữ liệu quan sát, thẩm định cảnh báo bằng \CTAV{} đến định vị nguyên nhân gốc có bằng chứng.

% -----------------------------------------------------------------------------
\section{Kết chương}
\label{sec:ch3-ket-chuong}

Mục này tổng hợp các nghiên cứu học thuật ở Mục~\ref{sec:ch3-nghien-cuu} và các hệ thống thực tế ở Mục~\ref{sec:ch3-he-thong} trên cùng một hệ tiêu chí, qua đó chỉ rõ bốn khoảng trống nghiên cứu và định vị đóng góp của đề tài.

\subsection{Bảng so sánh tổng hợp}
\label{sec:ch3-bang-so-sanh}

Bảng~\ref{tab:ch3-so-sanh} đối chiếu các nhóm công trình tiêu biểu và các hệ thống hiện có theo năm tiêu chí bám sát trực tiếp mục tiêu của đề tài: (1) khai thác đồ thị topology dịch vụ, (2) sử dụng LLM để thẩm định hoặc suy luận sự cố, (3) hiệu chỉnh quyết định theo nguồn gốc bằng chứng, (4) thực hiện phân tích nguyên nhân gốc (RCA), và (5) có đánh giá định lượng báo động giả (FP/FPR) trên quy trình mở tái lập được. Trong bảng, \emph{Có} biểu thị tiêu chí được đáp ứng trực tiếp; \emph{Một phần} biểu thị tiêu chí chỉ được đáp ứng gián tiếp hoặc ở phạm vi hẹp (chẳng hạn chỉ có luật ức chế tĩnh, chỉ học đồ thị cảm biến không phải topology dịch vụ, hoặc hệ thống thương mại đóng nguồn không công bố kiểm định FP độc lập); \emph{Không} biểu thị tiêu chí nằm ngoài phạm vi của công trình hoặc công cụ.

\begin{table}[htbp]
  \centering
  \caption[So sánh công trình liên quan, hệ thống hiện có và đề xuất]{So sánh các nhóm nghiên cứu liên quan, hệ thống hiện có và đề xuất của đề tài.}
  \label{tab:ch3-so-sanh}
  \footnotesize
  \renewcommand{\arraystretch}{1.15}
  \setlength{\tabcolsep}{3.5pt}
  \begin{tabularx}{\linewidth}{@{}Y >{\centering\arraybackslash}p{1.45cm} >{\centering\arraybackslash}p{1.45cm} >{\centering\arraybackslash}p{1.6cm} >{\centering\arraybackslash}p{1.2cm} >{\centering\arraybackslash}p{1.65cm}@{}}
    \toprule
    \textbf{Công trình / Hệ thống} &
    \textbf{Topology} &
    \textbf{LLM judge} &
    \textbf{Hiệu chỉnh BC} &
    \textbf{RCA} &
    \textbf{Đánh giá FP \& tái lập} \\
    \midrule
    Isolation Forest~\cite{liu2008iforest}, RRCF~\cite{guha2016rrcf} &
      Không & Không & Không & Không & Một phần \\
    MTAD-GAT~\cite{zhao2020mtadgat}, GDN~\cite{deng2021gdn} &
      Một phần & Không & Không & Không & Một phần \\
    LogParser-LLM~\cite{zhong2024logparser}, LogLLM~\cite{guan2024logllm} &
      Không & Có & Không & Không & Một phần \\
    RCACopilot~\cite{chen2024rcacopilot}, RCAgent~\cite{wang2024rcagent}, Roy~\cite{roy2024exploring} &
      Không & Có & Không & Có & Không \\
    RAAD-LLM~\cite{russellgilbert2025raad}, ARGOS~\cite{gu2025argos}, CALM~\cite{devireddy2025calm} &
      Không & Có & Không & Không & Một phần \\
    MonitorRank~\cite{kim2013monitorrank}, MicroRCA~\cite{wu2020microrca}, CIRCA~\cite{li2022circa}, RCD~\cite{ikram2022rcd} &
      Có & Không & Không & Có & Không \\
    Nezha~\cite{yu2023nezha}, BARO~\cite{pham2024baro}, RCAEval~\cite{pham2025rcaeval} &
      Có & Không & Không & Có & Một phần \\
    GraphRAG~\cite{edge2024graphrag} &
      Một phần & Có & Không & Không & Không \\
    Truy xuất topology cho thẩm định LLM~\cite{author2026icaetm} &
      Có & Có & Không & Không & Có \\
    FairEval~\cite{wang2024fair}, Xiong~\cite{xiong2024uncertainty}, ALCE~\cite{gao2023alce} &
      Không & Có & Một phần & Không & Một phần \\
    Alert Storm~\cite{zhao2020alertstorm}, Kuang~\cite{kuang2024knowledge} &
      Một phần & Không & Không & Không & Một phần \\
    \midrule
    Datadog (Event Mgmt + Bits AI)~\cite{datadog2026docs} &
      Có & Có & Không & Có & Không \\
    Dynatrace (Davis AI + Smartscape)~\cite{dynatrace2026docs} &
      Có & Không & Không & Có & Không \\
    PagerDuty AIOps~\cite{pagerduty2026docs} &
      Một phần & Không & Không & Một phần & Không \\
    Prometheus Alertmanager~\cite{prometheus2026alertmanager} &
      Không & Không & Không & Không & Một phần \\
    Keep~\cite{keep2026docs} &
      Có & Một phần & Không & Không & Một phần \\
    HolmesGPT~\cite{holmesgpt2026docs} &
      Một phần & Có & Không & Có & Không \\
    \midrule
    \textbf{Đề tài này (\CTAV{} + tác tử RCA)} &
      \textbf{Có} & \textbf{Có} & \textbf{Có} & \textbf{Có} & \textbf{Có} \\
    \bottomrule
  \end{tabularx}
\end{table}

Bảng~\ref{tab:ch3-so-sanh} cho thấy bức tranh phân mảnh rõ rệt giữa các hướng hiện có. Nhóm phát hiện bất thường chuỗi thời gian và kiến trúc giám sát lai chưa khai thác cấu trúc đồ thị dịch vụ; nhóm RCA dựa trên đồ thị và nhân quả khai thác tốt topology nhưng không dùng LLM thẩm định luồng cảnh báo thô để kìm báo động giả; nhóm tác tử RCA dùng LLM bỏ qua khâu lọc cảnh báo giả có kiểm soát ngân sách; còn nhóm hiệu chỉnh \emph{LLM-as-a-Judge} mới xét bài toán ngôn ngữ tổng quát chứ chưa hiệu chỉnh theo nguồn gốc bằng chứng trên đồ thị phụ thuộc. Các hệ thống công nghiệp hoặc đóng nguồn không thể tái lập kiểm định FP, hoặc tách rời tầng quản lý cảnh báo với tầng điều tra LLM.

\subsection{Khoảng trống nghiên cứu và định vị đề xuất}
\label{sec:ch3-khoang-trong}

Từ kết quả khảo sát và Bảng~\ref{tab:ch3-so-sanh}, đề tài xác định bốn khoảng trống nghiên cứu cụ thể:

\begin{enumerate}
  \item \textbf{Khoảng trống G1 — Khoảng cách giữa phát hiện tốc độ cao và thẩm định chính xác dưới ngân sách LLM hữu hạn.} Bộ phát hiện metric không giám sát chạy nhanh nhưng dễ tạo báo động giả hoặc bỏ sót lỗi ngữ nghĩa; ngược lại, tác tử LLM gọi nhiều công cụ trực tiếp quá tốn kém để chạy trên mọi cửa sổ thời gian. Cần một cơ chế kết nối có kiểm soát chặt số lượt gọi LLM mà vẫn bảo đảm chất lượng thẩm định.
  \item \textbf{Khoảng trống G2 — Đánh đổi chưa được kiểm soát giữa thu hồi theo topology và báo động giả do thiên lệch xác nhận.} Mở rộng truy xuất bằng chứng sang láng giềng trên đồ thị phụ thuộc giúp nhận diện các lỗi lan truyền mà truy xuất từ khóa cục bộ bỏ lỡ, nhưng đồng thời khuếch đại xu hướng xác nhận nhầm của \emph{LLM-as-a-Judge} trên các cửa sổ bình thường nếu không có chính sách lọc và neo bằng chứng phù hợp.
  \item \textbf{Khoảng trống G3 — Thiếu cơ chế hiệu chỉnh bộ thẩm định cảnh báo theo nguồn gốc bằng chứng và thiếu kiểm chứng loại trừ có đăng ký giả thuyết trước.} Các phương pháp hiệu chỉnh LLM hiện hành chưa mô hình hóa sự khác biệt giữa bằng chứng nội bộ và bằng chứng láng giềng hay tham chiếu ngưỡng nền từ tập phát triển. Đồng thời, trong thực nghiệm AIOps, nhiều thiết kế phức tạp được gộp chung vào một hệ thống mà không tách kiểm định từng thành phần trên tập kiểm tra giữ kín, dễ che giấu những thành phần không mang lại lợi ích thực tế.
  \item \textbf{Khoảng trống G4 — Sự đứt gãy giữa thẩm định cảnh báo đã kiểm soát FP với phân tích nguyên nhân gốc có trích dẫn bằng chứng trên một nền tảng mở.} Phần lớn công trình hoặc dừng ở bước phát hiện/thẩm định cảnh báo, hoặc bắt đầu RCA từ danh sách sự cố đã xác nhận sẵn, và hiếm khi kết hợp cả hai trên một nền tảng mã nguồn mở có bộ nhớ đệm lời gọi LLM để tái lập thực nghiệm và trình diễn trực quan.
\end{enumerate}

Để lấp đầy bốn khoảng trống trên, đề tài đề xuất \textbf{\CTAV{}} (\emph{Calibrated Topology-Aware Alert Validation}) và tích hợp \CTAV{} làm tầng nền cho \textbf{tác tử phân tích nguyên nhân gốc} trên một nền tảng thống nhất. Cụ thể, trong không gian thiết kế của \CTAV{}, đề tài xây dựng và kiểm định tách biệt từng thành phần theo quy trình đăng ký giả thuyết trước: thành phần \textbf{C1} (truy xuất lai BM25--topology, đang dùng trong cấu hình cuối) kết hợp với thành phần \textbf{C3} (neo bằng chứng theo ngưỡng nền đo trên tập phát triển \splitname{dev}, đang dùng trong cấu hình cuối) đóng vai trò lõi giúp duy trì lợi ích nhận diện của topology đồng thời kìm mạnh báo động giả. Hai cơ chế bổ sung trong thiết kế ban đầu là \textbf{C2} (ngưỡng quyết định bất đối xứng theo nguồn bằng chứng, đã loại sau kiểm định) và \textbf{C4} (định tuyến hai tầng thích nghi, đã loại sau kiểm định) được đánh giá minh bạch trên tập \splitname{val} và loại khỏi cấu hình cuối vì không tạo lợi ích bổ sung so với cặp C1+C3 tinh gọn; tương tự, cơ chế \textbf{C5} (\code{Rescue Score}, chỉ đánh giá ở vai trò phân tích loại trừ) được khảo sát độc lập thay vì mặc định bật trong hệ thống chính. Phía sau \CTAV{}, tác tử RCA gom các cảnh báo đã thẩm định thành sự cố, lần ngược đường lan truyền trên đồ thị phụ thuộc dịch vụ và sinh báo cáo chẩn đoán có cấu trúc kèm trích dẫn bằng chứng. Toàn bộ thiết kế kiến trúc, thuật toán chi tiết và mô hình dữ liệu của hệ thống được trình bày trong Chương~\ref{ch:de-xuat}.

%        ...
%        \bibliographystyle{IEEEtran}
%        \bibliography{refs}
%
% Cac nhan lien chuong duoc tham chieu tu Chuong 3:
%   Chuong: ch:mo-dau, ch:nen-tang, ch:de-xuat, ch:thuc-nghiem, ch:platform
%   Chuong 2: sec:ch2-du-lieu, sec:ch2-truy-xuat, sec:ch2-judge, sec:ch2-rca,
%             sec:ch2-dang-ky
%   Chuong 4: sec:ch4-ctav, sec:ch4-c1, sec:ch4-policy, sec:ch4-c4, sec:ch4-rca
%   Chuong 5: sec:ch5-hypotheses, sec:ch5-c5, sec:ch5-budget
% =============================================================================

\chapter{Công trình liên quan}
\label{ch:cong-trinh-lien-quan}

Chương này hệ thống hóa các hướng nghiên cứu và các hệ thống vận hành thực tế liên quan trực tiếp tới bài toán thẩm định cảnh báo có nhận biết topology và phân tích nguyên nhân gốc trên kiến trúc microservice. Nội dung được tổ chức xoay quanh hai câu hỏi: cộng đồng học thuật và công nghiệp đã giải quyết những khía cạnh nào của chuỗi xử lý từ phát hiện bất thường đến định vị nguyên nhân gốc, và còn những khoảng trống nào mà đề xuất của đề tài nhắm tới. Mục~\ref{sec:ch3-nghien-cuu} khảo sát bảy nhánh nghiên cứu học thuật từ phát hiện bất thường chuỗi thời gian, ứng dụng mô hình ngôn ngữ lớn trong phân tích log và nguyên nhân gốc, kiến trúc giám sát lai, phân tích nguyên nhân gốc dựa trên đồ thị, truy xuất dựa trên đồ thị, hiệu chỉnh độ tin cậy của \emph{LLM-as-a-Judge}, đến gom nhóm và phân loại cảnh báo. Mục~\ref{sec:ch3-he-thong} đối chiếu các nền tảng AIOps thương mại và công cụ mã nguồn mở hiện hành dựa trên tài liệu chính thức. Cuối cùng, Mục~\ref{sec:ch3-ket-chuong} tổng hợp các hướng tiếp cận trong Bảng~\ref{tab:ch3-so-sanh}, rút ra bốn khoảng trống nghiên cứu và định vị rõ phạm vi đóng góp của \CTAV{} cùng tác tử phân tích nguyên nhân gốc.

% -----------------------------------------------------------------------------
\section{Các nghiên cứu liên quan}
\label{sec:ch3-nghien-cuu}

Các công trình liên quan đến đề tài nằm ở giao điểm giữa giám sát chuỗi thời gian, phân tích dữ liệu quan sát đa phương thức, truy xuất bằng chứng theo đồ thị và hiệu chỉnh mô hình ngôn ngữ lớn đóng vai trò bộ thẩm định. Bảy tiểu mục dưới đây phân tích từng nhánh theo cùng một trục: cách tiếp cận của công trình, hạn chế đặt trong bài toán của đề tài và điểm khác biệt của đề tài.

\subsection{Phát hiện bất thường trên chuỗi thời gian metric}
\label{sec:ch3-tsad}

Hướng phát hiện bất thường trên chuỗi thời gian không giám sát hình thành nền tảng cho tầng lọc cảnh báo đầu vào. Liu và cộng sự~\cite{liu2008iforest} đề xuất Isolation Forest, một phương pháp dựa trên tập cây phân tách ngẫu nhiên tận dụng nhận xét rằng điểm bất thường ít và khác biệt nên cần ít phép chia để cô lập hơn điểm bình thường; mô hình chạy nhanh, không cần nhãn và tiêu thụ ít tài nguyên. Mở rộng sang dữ liệu luồng có phân bố biến đổi theo thời gian, Guha và cộng sự~\cite{guha2016rrcf} xây dựng Robust Random Cut Forest (RRCF) cho phép cập nhật cấu trúc cây trực tuyến khi điểm mới đến và đo độ bất thường bằng mức thay đổi độ phức tạp cây khi chèn hoặc xóa điểm. Hạn chế chung của nhóm cây phân tách là chỉ nhìn từng vector metric độc lập theo cửa sổ, nên không phân biệt được dao động tải bình thường với sự cố thực sự nếu không đối chiếu ngữ cảnh văn bản từ log hoặc quan hệ gọi giữa các service.

Để nắm bắt tương quan đa biến và phụ thuộc giữa các cảm biến hay metric, các công trình học sâu dựa trên đồ thị đã được đề xuất. Zhao và cộng sự~\cite{zhao2020mtadgat} giới thiệu MTAD-GAT với hai lớp chú ý đồ thị song song: một lớp học tương quan giữa các đặc trưng tại cùng thời điểm và một lớp học phụ thuộc dọc trục thời gian, kết hợp đồng thời mục tiêu dự báo và tái tạo để tính điểm bất thường. Deng và Hooi~\cite{deng2021gdn} đề xuất Graph Deviation Network (GDN), tự học đồ thị quan hệ giữa các cảm biến qua vector nhúng và dùng mạng chú ý đồ thị để dự báo giá trị tiếp theo, từ đó giải thích bất thường qua độ lệch của các cạnh dự báo. Tuy đạt độ chính xác cao hơn trên nhiều tập chuỗi thời gian đa biến, nhóm học sâu đồ thị đòi hỏi huấn luyện lại khi topology hoặc tập metric thay đổi, cần tài nguyên tính toán lớn hơn đáng kể và vẫn chỉ khai thác tín hiệu số mà bỏ qua nội dung log hay trace vốn chứa thông tin tường minh về lỗi ứng dụng.

Khác với các công trình tập trung thay thế bộ phát hiện metric bằng kiến trúc học sâu ngày càng phức tạp, đề tài này giữ tầng phát hiện nhẹ bằng Isolation Forest để bảo đảm tốc độ và khả năng chạy lại tức thì trên mọi cấu hình, rồi chuyển trọng tâm sang tầng thẩm định bằng bằng chứng đa phương thức phía sau. Cách phân công này cho phép tận dụng tốc độ lọc sơ bộ của detector nhẹ trong khi vẫn kiểm soát được báo động giả nhờ đối chiếu log và đồ thị phụ thuộc service.

\subsection{Mô hình ngôn ngữ lớn trong phân tích log và nguyên nhân gốc}
\label{sec:ch3-llm-log-rca}

Sự xuất hiện của các mô hình ngôn ngữ lớn (LLM) mở ra hướng xử lý trực tiếp dữ liệu log phi cấu trúc và suy luận sự cố bằng ngôn ngữ tự nhiên. Ở bài toán phân tách mẫu log (\emph{log parsing}), Zhong và cộng sự~\cite{zhong2024logparser} đề xuất LogParser-LLM, kết hợp cây tiền tố với truy vấn LLM chọn lọc để trích xuất mẫu log mà không cần huấn luyện lại mô hình khi định dạng log thay đổi. Ở bài toán phát hiện bất thường trên chuỗi log, Guan và cộng sự~\cite{guan2024logllm} giới thiệu LogLLM, phối hợp bộ mã hóa ngữ nghĩa dạng BERT cho từng dòng log với bộ giải mã LLM nhằm mô hình hóa chuỗi log mà không cần bước tách mẫu thủ công. Hai công trình này giải quyết tốt khâu biểu diễn và phân tích log đơn lẻ, nhưng chỉ làm việc trên tín hiệu log tách rời mà không tích hợp luồng cảnh báo từ metric cũng như cấu trúc phụ thuộc giữa các service.

Ở bài toán phân tích nguyên nhân gốc cho sự cố đám mây, nhiều công trình sử dụng LLM trực tiếp trên hồ sơ sự cố hoặc dưới dạng tác tử gọi công cụ. Ahmed và cộng sự~\cite{ahmed2023recommending} thực hiện đánh giá quy mô lớn trên hơn 40.000 sự cố vận hành thực tế tại Microsoft để khuyến nghị nguyên nhân gốc và bước khắc phục bằng mô hình dòng GPT, cho thấy LLM có ích trong việc tóm tắt tri thức sự cố nhưng nhạy cảm với chất lượng mô tả đầu vào. Chen và cộng sự~\cite{chen2024rcacopilot} xây dựng RCACopilot, một hệ thống định tuyến sự cố vào quy trình thu thập dữ liệu chẩn đoán tương ứng rồi để LLM dự đoán nhóm nguyên nhân gốc kèm giải thích bằng ngôn ngữ tự nhiên. Tiến thêm bước tự chủ, Wang và cộng sự~\cite{wang2024rcagent} đề xuất RCAgent sử dụng tác tử LLM có hỗ trợ công cụ truy vấn log và tổng hợp tự nhất quán (\emph{self-consistency}) trong môi trường đám mây, trong khi Roy và cộng sự~\cite{roy2024exploring} khảo sát tác tử theo vòng lặp ReAct trang bị công cụ truy xuất log, metric và cơ sở dữ liệu lịch sử để điều tra sự cố.

Điểm hạn chế chung của nhóm công trình LLM cho RCA nêu trên là phần lớn giả định sự cố đầu vào đã được xác nhận hoặc đã mở phiếu sự cố sẵn, nên không đánh giá trực tiếp bài toán thẩm định luồng cảnh báo thô để chặn báo động giả trước khi gọi tác tử chẩn đoán tốn kém. Ngoài ra, các hệ thống đánh giá trên dữ liệu sự cố nội bộ doanh nghiệp khó tái lập độc lập bên ngoài. Đề tài này đặt tầng thẩm định cảnh báo có kiểm soát ngân sách ngay trước tác tử RCA và thực nghiệm hoàn toàn trên bộ dữ liệu công khai có thể tái lập.

\subsection{Kiến trúc giám sát lai và LLM-as-a-Judge}
\label{sec:ch3-hybrid-judge}

Để dung hòa chi phí gọi LLM với yêu cầu giám sát liên tục, nhiều công trình kết hợp mô hình nhẹ hoặc cơ chế truy xuất tăng cường sinh (\emph{Retrieval-Augmented Generation}, RAG)~\cite{lewis2020rag} với bộ phân xử LLM ở tầng sau. Về nguyên lý phân tầng chi phí, nền tảng lý thuyết bắt nguồn từ cơ chế đánh đổi sai số và quyền từ chối (\emph{reject option}) của Chow~\cite{chow1970reject} cùng mô hình chuỗi tầng gọi LLM tiết kiệm chi phí FrugalGPT của Chen và cộng sự~\cite{chen2024frugalgpt}, trong đó các mẫu rõ ràng được xử lý bởi mô hình rẻ hơn và chỉ vùng bất định mới chuyển lên mô hình đắt hơn.

Trong miền phát hiện bất thường, một loạt công trình gần đây đưa kiến trúc lai vào thực nghiệm. Russell-Gilbert và cộng sự~\cite{russellgilbert2025raad} đề xuất RAAD-LLM, tích hợp RAG vào LLM để thích nghi phát hiện bất thường chuỗi thời gian khi bối cảnh hệ thống thay đổi mà không cần tinh chỉnh trọng số. Gu và cộng sự~\cite{gu2025argos} giới thiệu ARGOS, một khuôn khổ tác tử cho phép LLM tự sinh và tinh chỉnh luật phát hiện bất thường trên chuỗi thời gian nhằm giảm phụ thuộc vào truy vấn trực tiếp trên từng mẫu. Devireddy và Huang~\cite{devireddy2025calm} đề xuất CALM, điều phối liên tục giữa phát hiện thời gian thực và cơ chế can thiệp thích nghi của LLM. Ở góc độ đánh giá năng lực thẩm định, Sivalingam và cộng sự~\cite{sivalingam2026tsquery} xây dựng TSQueryBench để kiểm tra khả năng dùng \emph{LLM-as-a-Judge} đánh giá lời giải thích bất thường chuỗi thời gian.

Hạn chế chung của các kiến trúc lai này khi áp vào microservice là chúng chủ yếu vận hành trên chuỗi thời gian đơn lẻ hoặc truy xuất ngữ cảnh văn bản phẳng, chưa kết nối tường minh với đồ thị phụ thuộc giữa các service sinh từ trace. Quan trọng hơn, các nghiên cứu này chưa khảo sát hệ quả ngược chiều khi mở rộng bằng chứng cho bộ thẩm định LLM: việc cấp thêm ngữ cảnh từ láng giềng có thể làm tăng tỷ lệ báo động giả do thiên lệch xác nhận. Đề tài này tập trung giải quyết trực tiếp bài toán kiểm soát ngân sách và kìm báo động giả trong kiến trúc lai detector nhẹ kết hợp \emph{LLM-as-a-Judge} cho hệ microservice.

\subsection{Phân tích nguyên nhân gốc dựa trên topology, đồ thị và suy luận nhân quả}
\label{sec:ch3-rca-graph}

Định vị nguyên nhân gốc dựa trên đồ thị phụ thuộc và suy luận nhân quả là nhánh phát triển lâu đời nhất trong giám sát hệ phân tán. Kim và cộng sự~\cite{kim2013monitorrank} đề xuất MonitorRank, thực hiện bước đi ngẫu nhiên (\emph{random walk}) có trọng số tương quan metric trên đồ thị gọi dịch vụ để xếp hạng service nghi vấn trong kiến trúc hướng dịch vụ. Chen và cộng sự~\cite{chen2014causeinfer} xây dựng CauseInfer với đồ thị nhân quả hai tầng (tầng dịch vụ và tầng metric) bằng thuật toán kiểm định độc lập có điều kiện PC để lần ngược nút gốc gây suy giảm hiệu năng. Trên kiến trúc microservice, Wu và cộng sự~\cite{wu2020microrca} giới thiệu MicroRCA, gắn dữ liệu sử dụng tài nguyên máy chủ vào đồ thị dịch vụ bất thường và áp dụng Personalized PageRank để định vị service gây lỗi hiệu năng.

Các công trình gần đây tiếp cận bài toán bằng mô hình nhân quả chặt chẽ hơn và dữ liệu đa phương thức. Li và cộng sự~\cite{li2022circa} đề xuất CIRCA, mô hình hóa lỗi như một can thiệp (\emph{intervention}) trên mạng Bayes nhân quả xây từ kiến trúc hệ thống và dùng điều kiện thay đổi phân bố có điều kiện để chỉ ra chỉ số gốc. Ikram và cộng sự~\cite{ikram2022rcd} đề xuất RCD, xem sự cố như can thiệp cục bộ và khám phá đồ thị nhân quả theo phân cấp để tránh chi phí tính toán lớn của thuật toán PC trên toàn bộ hệ thống. Kết hợp cả ba trụ cột dữ liệu quan sát, Yu và cộng sự~\cite{yu2023nezha} đề xuất Nezha, chuyển đổi metric, log và trace thành biểu diễn sự kiện đồng nhất và so sánh mẫu phân kỳ giữa trạng thái bình thường với trạng thái lỗi để định vị nguyên nhân gốc ở mức chi tiết, đồng thời công bố bộ dữ liệu microservice đa phương thức làm chuẩn đánh giá. Đáng lưu ý về tính bền vững của chuỗi xử lý, Pham và cộng sự~\cite{pham2024baro} đề xuất BARO bằng phát hiện điểm đổi Bayes đa biến kết hợp kiểm định phi tham số, đồng thời chỉ ra rằng sai số ở khâu phát hiện bất thường đầu vào có thể lan sang làm giảm mạnh độ chính xác định vị gốc; tiếp đó nhóm tác giả công bố bộ chuẩn RCAEval~\cite{pham2025rcaeval} tổng hợp 735 ca lỗi đa phương thức trên nhiều hệ microservice mã nguồn mở để thống nhất việc so sánh.

Hạn chế của nhóm phương pháp đồ thị và nhân quả thuần túy gồm hai điểm. Thứ nhất, phần lớn mô hình chỉ xuất ra danh sách xếp hạng service hoặc metric gốc dưới dạng chỉ số toán học, thiếu báo cáo tường minh giải thích cơ chế lan truyền bằng chứng từ log và trace cho kỹ sư vận hành. Thứ hai, trừ những phân tích gần đây như BARO, đa số nghiên cứu RCA giả định cửa sổ thời gian sự cố đã được xác định sạch từ trước, tách rời khỏi bài toán lọc cảnh báo giả trực tuyến. Đề tài này nối trực tiếp hai khâu: dùng \CTAV{} để làm sạch luồng cảnh báo trước khi gom sự cố, rồi dùng tác tử RCA lần ngược chuỗi lan truyền trên đồ thị phụ thuộc và sinh báo cáo có cấu trúc kèm trích dẫn bằng chứng.

\subsection{Truy xuất dựa trên đồ thị và GraphRAG cho vận hành hệ thống}
\label{sec:ch3-graphrag}

Trong lĩnh vực xử lý ngôn ngữ tự nhiên, Edge và cộng sự~\cite{edge2024graphrag} đề xuất GraphRAG nhằm khắc phục hạn chế của truy xuất văn bản phẳng khi cần tổng hợp thông tin phân tán trên nhiều tài liệu. Phương pháp xây dựng đồ thị tri thức thực thể--quan hệ từ kho văn bản, phân cụm cộng đồng đồ thị và sinh tóm tắt phân cấp để trả lời các truy vấn mang tính toàn cục. Tuy mở ra hướng truy xuất có cấu trúc, GraphRAG thiết kế cho kho tài liệu tĩnh với đồ thị trích xuất bằng LLM, nên vừa tốn chi phí lập chỉ mục vừa khác về bản chất so với đồ thị phụ thuộc động của hệ microservice được đo trực tiếp từ quan hệ gọi \code{SpanID}--\code{ParentID} trong trace.

Khi áp dụng tư tưởng truy xuất theo cấu trúc đồ thị vào bài toán thẩm định cảnh báo microservice, một nghiên cứu tiền đề về truy xuất bằng chứng có nhận biết topology~\cite{author2026icaetm} đã đề xuất kết hợp điểm từ khóa BM25 với độ gần về thời gian và khoảng cách bước nhảy trên đồ thị phụ thuộc service để cung cấp thêm bằng chứng từ các service láng giềng cho bộ thẩm định LLM (chi tiết công thức hợp nhất đã trình bày ở Mục~\ref{sec:ch2-truy-xuat}). Kết quả định tính của hướng tiếp cận này cho thấy một đánh đổi rõ rệt: việc mở rộng ngữ cảnh sang các service có liên kết gọi giúp cải thiện khả năng nhận diện các sự cố lan truyền mà bằng chứng cục bộ không đủ, nhưng khi đưa trực tiếp vào \emph{LLM-as-a-Judge} mà không có cơ chế kiểm soát chính sách, số báo động giả tăng lên mạnh do bộ thẩm định dễ diễn giải nhầm log nền từ láng giềng thành dấu hiệu bất thường. Đề tài này xuất phát trực tiếp từ đánh đổi đó để xây dựng cơ chế hiệu chỉnh theo nguồn gốc bằng chứng và kiểm định nghiêm ngặt hiệu quả kìm báo động giả trên dữ liệu ngoài tập tinh chỉnh.

\subsection{Hiệu chỉnh độ tin cậy, thiên lệch của LLM-as-a-Judge và gắn bằng chứng}
\label{sec:ch3-calibration}

Dù \emph{LLM-as-a-Judge} đạt mức đồng thuận cao với người đánh giá trên nhiều tác vụ tổng quát~\cite{zheng2023judging}, hàng loạt nghiên cứu đã chứng minh mô hình ngôn ngữ mắc nhiều thiên lệch có hệ thống và thường tự tin thái quá. Wang và cộng sự~\cite{wang2024fair} chỉ ra rằng bộ thẩm định LLM chịu thiên lệch vị trí nghiêm trọng khi so sánh cặp câu trả lời, dẫn đến thứ hạng dễ bị thao túng chỉ bằng cách đảo thứ tự trình bày, và đề xuất khung FairEval hiệu chỉnh bằng cách buộc sinh bằng chứng trước khi chấm điểm kết hợp trung bình hóa vị trí. Sharma và cộng sự~\cite{sharma2024sycophancy} phân tích hiện tượng chiều lòng (\emph{sycophancy}), cho thấy mô hình đã qua tinh chỉnh theo phản hồi con người có xu hướng đồng thuận với giả định ngầm trong câu hỏi ngay cả khi giả định đó sai — một cơ chế tương đồng trực tiếp với thiên lệch xác nhận khi bộ thẩm định nhận thông báo rằng ``cảnh báo đang kích hoạt''.

Về hiệu chỉnh xác suất và ước lượng độ bất định, Guo và cộng sự~\cite{guo2017calibration} chứng minh các mạng nơ-ron hiện đại tuy chính xác nhưng hiệu chỉnh kém và đề xuất phương pháp hiệu chỉnh hậu xử lý bằng tỉ lệ nhiệt độ (\emph{temperature scaling}). Trên LLM, Kadavath và cộng sự~\cite{kadavath2022mostly} khảo sát khả năng mô hình tự đánh giá xác suất đúng của phát biểu do chính nó sinh ra, cho thấy tín hiệu nội tại có thể hiệu chỉnh được trong câu hỏi trắc nghiệm nhưng khó khái quát khi chuyển phân bố tác vụ. Gần hơn với bài toán sử dụng API thương mại không truy cập được xác suất token thô, Xiong và cộng sự~\cite{xiong2024uncertainty} đánh giá thực nghiệm độ tin cậy bằng lời (\emph{verbalized confidence}) trên nhiều dòng LLM và chỉ ra mô hình thường gán điểm tin cậy quá cao, tập trung dày đặc ở vùng sát $1{,}0$. Song song đó, để hạn chế suy diễn vô căn cứ trong văn bản sinh ra, Gao và cộng sự~\cite{gao2023alce} đề xuất chuẩn đánh giá ALCE đo đồng thời tính đúng đắn và chất lượng trích dẫn nguồn bằng chứng trong câu trả lời của LLM.

Khoảng trống còn lại trong nhánh nghiên cứu này là các kỹ thuật hiệu chỉnh hiện có chủ yếu thiết kế cho bài toán chấm điểm cặp câu trả lời, hỏi--đáp mở hoặc sinh văn bản có trích dẫn tổng quát. Chưa có cơ chế nào hiệu chỉnh trực tiếp quyết định của \emph{LLM-as-a-Judge} trong bài toán thẩm định cảnh báo microservice theo \emph{nguồn gốc bằng chứng} (nội bộ hay láng giềng trên đồ thị phụ thuộc) và theo \emph{mức độ bất thường nền đo trên tập phát triển}. Đề tài này cụ thể hóa ý tưởng hiệu chỉnh bằng chính sách neo bằng chứng và kiểm định tách biệt từng thành phần thiết kế (ngưỡng bất đối xứng theo nguồn bằng chứng, neo ngưỡng nền trên tập phát triển và định tuyến hai tầng) theo quy trình đăng ký giả thuyết trước ở Mục~\ref{sec:ch2-dang-ky}.

\subsection{Gom nhóm, phân loại và cứu cảnh báo trong vận hành}
\label{sec:ch3-alert-triage}

Trong vận hành hệ thống trực tuyến quy mô lớn, xử lý bão cảnh báo và định tuyến sự cố là bài toán sống còn để giảm tải cho kỹ sư trực. Zhao và cộng sự~\cite{zhao2020alertstorm} khảo sát hiện tượng bão cảnh báo trong hệ dịch vụ trực tuyến công nghiệp và đề xuất quy trình hai bước gồm phát hiện thời điểm bão và gom nhóm các cảnh báo tương quan về thời gian--ngữ nghĩa thành một cụm sự cố đại diện. Ở khâu định tuyến sự cố, Chen và cộng sự~\cite{chen2019empirical,chen2019continuous} phân tích nguyên nhân định tuyến nhầm sự cố tại các dịch vụ đám mây lớn và xây dựng phương pháp phân loại liên tục tận dụng thông tin cập nhật trong quá trình thảo luận của kỹ sư. Gần đây hơn, Kuang và cộng sự~\cite{kuang2024knowledge} đề xuất phương pháp gom cảnh báo có nhận biết tri thức cho hệ đám mây quy mô lớn, kết hợp đặc trưng văn bản với tri thức vận hành và topology để gộp cảnh báo cùng gốc tốt hơn các luật tĩnh.

Các công trình gom nhóm và phân loại này giải quyết tốt khâu \emph{sau khi cảnh báo đã phát sinh}, giúp nén số lượng thông báo gửi tới người trực, nhưng thường không thẩm định lại bản thân cảnh báo đầu vào có phải là báo động giả do bộ phát hiện metric gây ra hay không. Mặt khác, một câu hỏi ít được kiểm chứng nghiêm ngặt là liệu các cửa sổ bị bộ phát hiện chính bỏ sót có thể được ``cứu'' lại bằng điểm phụ trợ từ tín hiệu log/trace (cơ chế \code{Rescue Score}) mà không làm vỡ tỷ lệ báo động giả hay không. Đề tài này tách rõ hai vấn đề: sử dụng gom nhóm theo cửa sổ và topology ở đầu vào tác tử RCA để hợp nhất các cảnh báo đã được \CTAV{} thẩm định thành sự cố, đồng thời đánh giá cơ chế cứu cảnh báo như một phân tích loại trừ độc lập (Mục~\ref{sec:ch5-c5}) thay vì mặc định đưa vào hệ thống chính.

% -----------------------------------------------------------------------------
\section{Các hệ thống và công cụ hiện có}
\label{sec:ch3-he-thong}

Bên cạnh các công trình học thuật, nhiều nền tảng AIOps thương mại và công cụ mã nguồn mở đã được triển khai rộng rãi trong thực tế vận hành. Mục này mô tả tính năng cốt lõi của các hệ thống tiêu biểu dựa trên tài liệu chính thức của nhà phát triển, sau đó chỉ ra điểm khác biệt về mục tiêu và giới hạn khi đặt trong bối cảnh nghiên cứu của đề tài.

\subsection{Nền tảng AIOps thương mại}
\label{sec:ch3-thuong-mai}

\textbf{Datadog (Event Management và Bits AI).} Theo tài liệu chính thức của Datadog~\cite{datadog2026docs}, phân hệ Event Management thu nạp tập trung cảnh báo từ các bộ giám sát nội bộ cùng nguồn bên thứ ba, sau đó áp dụng tương quan cảnh báo tự động dựa trên topology dịch vụ, nhãn chia sẻ và mẫu thời gian để gom các cảnh báo lan truyền vào chung một sự vụ, đồng thời làm giàu ngữ cảnh về đơn vị sở hữu dịch vụ. Ở lớp điều tra sự cố, tác tử Bits AI (Bits Investigation) được kích hoạt tự động khi có cảnh báo nghiêm trọng, thực hiện kiểm định song song nhiều giả thuyết trên dữ liệu metric, log, trace và thay đổi cấu hình để lần theo chuỗi nhân quả và đề xuất hướng xử lý. Tuy cung cấp quy trình khép kín từ gom cảnh báo đến điều tra bằng AI, kiến trúc bên trong của cơ chế tương quan và điều tra đều đóng kín; nhà cung cấp không công bố phương pháp hiệu chỉnh độ tin cậy của bộ suy luận LLM theo nguồn bằng chứng hay đánh giá tách biệt tác động của topology lên tỷ lệ báo động giả trên bộ chuẩn công khai.

\textbf{Dynatrace (Davis AI và Smartscape).} Tài liệu kỹ thuật của Dynatrace~\cite{dynatrace2026docs} mô tả kiến trúc dựa trên hai trụ cột là bản đồ phụ thuộc thời gian thực Smartscape và động cơ phân tích nhân quả Davis AI. Smartscape tự động lập đồ thị phụ thuộc đầy đủ từ ứng dụng, service, tiến trình đến máy chủ và hạ tầng Kubernetes. Khi phát sinh bất thường ở mức chỉ số, độ trễ hoặc tỷ lệ lỗi, Davis AI gom các sự kiện liên quan có cùng nguyên nhân gốc thành một thực thể \emph{Problem} duy nhất và thực hiện phân tích cây lỗi tất định (\emph{deterministic fault-tree analysis}) trên đồ thị Smartscape để tách biệt thành phần gốc khỏi các thành phần chịu ảnh hưởng, hiển thị qua biểu đồ \emph{Visual Resolution Path}. Ưu điểm của Dynatrace là tính nhất quán cao nhờ mô hình nhân quả dựa trên agent thu thập sâu (\emph{OneAgent}), nhưng hệ thống gắn chặt với hạ tầng thu thập độc quyền, chi phí bản quyền lớn và không nhằm giải bài toán hiệu chỉnh quyết định của \emph{LLM-as-a-Judge} trên dữ liệu log phi cấu trúc lẫn ngân sách truy vấn LLM giới hạn.

\textbf{PagerDuty AIOps.} Theo tài liệu hỗ trợ chính thức của PagerDuty~\cite{pagerduty2026docs}, bộ công cụ PagerDuty AIOps tập trung vào ba chức năng: \emph{Intelligent Alert Grouping} dùng học máy nhóm các cảnh báo có liên quan theo mẫu lịch sử để giảm nhiễu thông báo; \emph{Probable Origin} phân tích chuỗi sự cố lan truyền trong lịch sử sáu tháng gần nhất để gợi ý tối đa ba dịch vụ hoặc sự cố có khả năng là điểm khởi phát cao nhất; và \emph{Event Orchestration} cung cấp bộ máy luật cho phép lọc, làm giàu và định tuyến sự kiện tự động ở mức toàn cục hoặc từng dịch vụ. Cách tiếp cận của PagerDuty giải quyết hiệu quả bài toán quản trị phản ứng sự cố ở tầng thông báo, nhưng phụ thuộc nhiều vào lịch sử sự cố tích lũy theo tổ chức, không khai thác đồ thị gọi chi tiết từng cửa sổ trace để thẩm định cảnh báo từ tín hiệu thô và không công bố cơ chế kiểm soát báo động giả cho tác tử thẩm định LLM.

\subsection{Công cụ mã nguồn mở}
\label{sec:ch3-ma-nguon-mo}

\textbf{Prometheus Alertmanager.} Theo tài liệu chính thức của dự án Prometheus~\cite{prometheus2026alertmanager}, Alertmanager đảm nhận tầng xử lý thông báo sau khi các luật cảnh báo Prometheus kích hoạt, cung cấp bốn cơ chế nền tảng: cây định tuyến theo nhãn (\code{route}), gom nhóm cảnh báo cùng nhãn thành một thông báo (\code{group\_by}), ức chế cảnh báo đích khi cảnh báo nguồn cùng nhóm đang kích hoạt (\code{inhibit\_rules}) và tắt tiếng tạm thời theo khung bảo trì (\code{silences}). Alertmanager nhẹ, minh bạch và là tiêu chuẩn thực tế trong hệ sinh thái đám mây, nhưng mọi quan hệ ức chế và gom nhóm đều phải cấu hình tĩnh bằng tay trong tệp YAML; công cụ không tự suy diễn đồ thị phụ thuộc từ trace, không đọc nội dung log và không thực hiện thẩm định cảnh báo hay phân tích nguyên nhân gốc.

\textbf{Keep.} Keep là nền tảng AIOps và quản lý cảnh báo mã nguồn mở; theo tài liệu chính thức của dự án~\cite{keep2026docs}, hệ thống hỗ trợ khử trùng lặp cảnh báo bằng dấu vân tay (\emph{fingerprint deduplication}, gồm chế độ toàn phần và từng phần), cung cấp ba cơ chế tương quan cảnh báo (luật thủ công bằng biểu thức CEL, tương quan theo topology và tương quan hỗ trợ bởi AI) cùng quy trình tự động hóa khai báo bằng YAML để làm giàu dữ liệu và gọi công cụ bên ngoài. Keep khắc phục được tính đơn điệu của công cụ luật tĩnh và mở rộng kết nối tốt với nhiều nguồn giám sát, nhưng vẫn xoay quanh bài toán tổng hợp--điều phối cảnh báo và quy trình tự động hóa; tài liệu hiện hành chưa đề cập cơ chế hiệu chỉnh \emph{LLM-as-a-Judge} theo nguồn bằng chứng hay khung thực nghiệm tái lập đánh giá đánh đổi recall--FPR dưới ngân sách LLM cố định.

\textbf{HolmesGPT.} HolmesGPT là tác tử điều tra sự cố mã nguồn mở khởi nguồn từ Robusta và hiện là dự án CNCF Sandbox; theo tài liệu chính thức~\cite{holmesgpt2026docs}, công cụ vận hành theo vòng lặp ReAct chạy ngay trong cụm Kubernetes ở chế độ chỉ đọc và tuân thủ phân quyền RBAC. Khi nhận cảnh báo từ Alertmanager hoặc PagerDuty, HolmesGPT để LLM tự chọn và gọi các công cụ truy vấn dữ liệu quan sát (như Prometheus, Loki hoặc \code{kubectl}), đọc kết quả trả về qua nhiều bước rồi tổng hợp nhận định nguyên nhân gốc. HolmesGPT minh họa rõ tính khả thi của tác tử LLM mã nguồn mở trong vận hành thực tế, nhưng chủ yếu đảm nhận khâu điều tra sau khi cảnh báo đã gửi tới, tốn nhiều lượt gọi LLM cho mỗi cảnh báo theo vòng lặp mở, và chưa tích hợp tầng thẩm định hiệu chỉnh trước đó để chặn báo động giả hay khung đánh giá thực nghiệm thống kê trên bộ chuẩn công khai.

\subsection{Nhận xét về các hệ thống hiện có}
\label{sec:ch3-nhan-xet}

Qua đối chiếu ba nền tảng thương mại và ba công cụ mã nguồn mở, có thể rút ra hai nhận xét chính về hiện trạng công nghiệp. Thứ nhất, các nền tảng thương mại (Datadog, Dynatrace, PagerDuty) đã tích hợp sâu khả năng gom cảnh báo theo topology và từng bước đưa tác tử AI vào điều tra sự cố, nhưng tất cả đều là hệ thống hộp đen gắn với hệ sinh thái thu thập độc quyền; chúng không cho phép kiểm định độc lập tác động của từng thành phần truy xuất hay chính sách hiệu chỉnh lên tỷ lệ báo động giả và không phục vụ được mục tiêu tái lập thực nghiệm khoa học. Thứ hai, ở mảng mã nguồn mở, hai năng lực cốt lõi vẫn nằm tách rời: nhóm quản lý cảnh báo (Alertmanager, Keep) mạnh về khử trùng lặp, luật ức chế và quy trình gom nhóm nhưng thiếu bộ thẩm định dựa trên bằng chứng log--trace đã hiệu chỉnh, còn tác tử điều tra bằng LLM (HolmesGPT) lại chạy trực tiếp theo vòng lặp nhiều bước mà không có tầng lọc kiểm soát ngân sách và kìm báo động giả ở đầu vào.

Do đó, thay vì cạnh tranh về độ phủ kết nối hạ tầng sản xuất với các sản phẩm thương mại, nền tảng trong đề tài này được xây dựng như một kiến trúc mở gọn nhẹ phục vụ đồng thời hai mục đích: tái lập chặt chẽ các thí nghiệm so sánh và loại trừ trên bộ dữ liệu công khai, và trình diễn trọn vẹn luồng xử lý từ phát lại dữ liệu quan sát, thẩm định cảnh báo bằng \CTAV{} đến định vị nguyên nhân gốc có bằng chứng.

% -----------------------------------------------------------------------------
\section{Kết chương}
\label{sec:ch3-ket-chuong}

Mục này tổng hợp các nghiên cứu học thuật ở Mục~\ref{sec:ch3-nghien-cuu} và các hệ thống thực tế ở Mục~\ref{sec:ch3-he-thong} trên cùng một hệ tiêu chí, qua đó chỉ rõ bốn khoảng trống nghiên cứu và định vị đóng góp của đề tài.

\subsection{Bảng so sánh tổng hợp}
\label{sec:ch3-bang-so-sanh}

Bảng~\ref{tab:ch3-so-sanh} đối chiếu các nhóm công trình tiêu biểu và các hệ thống hiện có theo năm tiêu chí bám sát trực tiếp mục tiêu của đề tài: (1) khai thác đồ thị topology dịch vụ, (2) sử dụng LLM để thẩm định hoặc suy luận sự cố, (3) hiệu chỉnh quyết định theo nguồn gốc bằng chứng, (4) thực hiện phân tích nguyên nhân gốc (RCA), và (5) có đánh giá định lượng báo động giả (FP/FPR) trên quy trình mở tái lập được. Trong bảng, \emph{Có} biểu thị tiêu chí được đáp ứng trực tiếp; \emph{Một phần} biểu thị tiêu chí chỉ được đáp ứng gián tiếp hoặc ở phạm vi hẹp (chẳng hạn chỉ có luật ức chế tĩnh, chỉ học đồ thị cảm biến không phải topology dịch vụ, hoặc hệ thống thương mại đóng nguồn không công bố kiểm định FP độc lập); \emph{Không} biểu thị tiêu chí nằm ngoài phạm vi của công trình hoặc công cụ.

\begin{table}[htbp]
  \centering
  \caption[So sánh công trình liên quan, hệ thống hiện có và đề xuất]{So sánh các nhóm nghiên cứu liên quan, hệ thống hiện có và đề xuất của đề tài.}
  \label{tab:ch3-so-sanh}
  \footnotesize
  \renewcommand{\arraystretch}{1.15}
  \setlength{\tabcolsep}{3.5pt}
  \begin{tabularx}{\linewidth}{@{}Y >{\centering\arraybackslash}p{1.45cm} >{\centering\arraybackslash}p{1.45cm} >{\centering\arraybackslash}p{1.6cm} >{\centering\arraybackslash}p{1.2cm} >{\centering\arraybackslash}p{1.65cm}@{}}
    \toprule
    \textbf{Công trình / Hệ thống} &
    \textbf{Topology} &
    \textbf{LLM judge} &
    \textbf{Hiệu chỉnh BC} &
    \textbf{RCA} &
    \textbf{Đánh giá FP \& tái lập} \\
    \midrule
    Isolation Forest~\cite{liu2008iforest}, RRCF~\cite{guha2016rrcf} &
      Không & Không & Không & Không & Một phần \\
    MTAD-GAT~\cite{zhao2020mtadgat}, GDN~\cite{deng2021gdn} &
      Một phần & Không & Không & Không & Một phần \\
    LogParser-LLM~\cite{zhong2024logparser}, LogLLM~\cite{guan2024logllm} &
      Không & Có & Không & Không & Một phần \\
    RCACopilot~\cite{chen2024rcacopilot}, RCAgent~\cite{wang2024rcagent}, Roy~\cite{roy2024exploring} &
      Không & Có & Không & Có & Không \\
    RAAD-LLM~\cite{russellgilbert2025raad}, ARGOS~\cite{gu2025argos}, CALM~\cite{devireddy2025calm} &
      Không & Có & Không & Không & Một phần \\
    MonitorRank~\cite{kim2013monitorrank}, MicroRCA~\cite{wu2020microrca}, CIRCA~\cite{li2022circa}, RCD~\cite{ikram2022rcd} &
      Có & Không & Không & Có & Không \\
    Nezha~\cite{yu2023nezha}, BARO~\cite{pham2024baro}, RCAEval~\cite{pham2025rcaeval} &
      Có & Không & Không & Có & Một phần \\
    GraphRAG~\cite{edge2024graphrag} &
      Một phần & Có & Không & Không & Không \\
    Truy xuất topology cho thẩm định LLM~\cite{author2026icaetm} &
      Có & Có & Không & Không & Có \\
    FairEval~\cite{wang2024fair}, Xiong~\cite{xiong2024uncertainty}, ALCE~\cite{gao2023alce} &
      Không & Có & Một phần & Không & Một phần \\
    Alert Storm~\cite{zhao2020alertstorm}, Kuang~\cite{kuang2024knowledge} &
      Một phần & Không & Không & Không & Một phần \\
    \midrule
    Datadog (Event Mgmt + Bits AI)~\cite{datadog2026docs} &
      Có & Có & Không & Có & Không \\
    Dynatrace (Davis AI + Smartscape)~\cite{dynatrace2026docs} &
      Có & Không & Không & Có & Không \\
    PagerDuty AIOps~\cite{pagerduty2026docs} &
      Một phần & Không & Không & Một phần & Không \\
    Prometheus Alertmanager~\cite{prometheus2026alertmanager} &
      Không & Không & Không & Không & Một phần \\
    Keep~\cite{keep2026docs} &
      Có & Một phần & Không & Không & Một phần \\
    HolmesGPT~\cite{holmesgpt2026docs} &
      Một phần & Có & Không & Có & Không \\
    \midrule
    \textbf{Đề tài này (\CTAV{} + tác tử RCA)} &
      \textbf{Có} & \textbf{Có} & \textbf{Có} & \textbf{Có} & \textbf{Có} \\
    \bottomrule
  \end{tabularx}
\end{table}

Bảng~\ref{tab:ch3-so-sanh} cho thấy bức tranh phân mảnh rõ rệt giữa các hướng hiện có. Nhóm phát hiện bất thường chuỗi thời gian và kiến trúc giám sát lai chưa khai thác cấu trúc đồ thị dịch vụ; nhóm RCA dựa trên đồ thị và nhân quả khai thác tốt topology nhưng không dùng LLM thẩm định luồng cảnh báo thô để kìm báo động giả; nhóm tác tử RCA dùng LLM bỏ qua khâu lọc cảnh báo giả có kiểm soát ngân sách; còn nhóm hiệu chỉnh \emph{LLM-as-a-Judge} mới xét bài toán ngôn ngữ tổng quát chứ chưa hiệu chỉnh theo nguồn gốc bằng chứng trên đồ thị phụ thuộc. Các hệ thống công nghiệp hoặc đóng nguồn không thể tái lập kiểm định FP, hoặc tách rời tầng quản lý cảnh báo với tầng điều tra LLM.

\subsection{Khoảng trống nghiên cứu và định vị đề xuất}
\label{sec:ch3-khoang-trong}

Từ kết quả khảo sát và Bảng~\ref{tab:ch3-so-sanh}, đề tài xác định bốn khoảng trống nghiên cứu cụ thể:

\begin{enumerate}
  \item \textbf{Khoảng trống G1 — Khoảng cách giữa phát hiện tốc độ cao và thẩm định chính xác dưới ngân sách LLM hữu hạn.} Bộ phát hiện metric không giám sát chạy nhanh nhưng dễ tạo báo động giả hoặc bỏ sót lỗi ngữ nghĩa; ngược lại, tác tử LLM gọi nhiều công cụ trực tiếp quá tốn kém để chạy trên mọi cửa sổ thời gian. Cần một cơ chế kết nối có kiểm soát chặt số lượt gọi LLM mà vẫn bảo đảm chất lượng thẩm định.
  \item \textbf{Khoảng trống G2 — Đánh đổi chưa được kiểm soát giữa thu hồi theo topology và báo động giả do thiên lệch xác nhận.} Mở rộng truy xuất bằng chứng sang láng giềng trên đồ thị phụ thuộc giúp nhận diện các lỗi lan truyền mà truy xuất từ khóa cục bộ bỏ lỡ, nhưng đồng thời khuếch đại xu hướng xác nhận nhầm của \emph{LLM-as-a-Judge} trên các cửa sổ bình thường nếu không có chính sách lọc và neo bằng chứng phù hợp.
  \item \textbf{Khoảng trống G3 — Thiếu cơ chế hiệu chỉnh bộ thẩm định cảnh báo theo nguồn gốc bằng chứng và thiếu kiểm chứng loại trừ có đăng ký giả thuyết trước.} Các phương pháp hiệu chỉnh LLM hiện hành chưa mô hình hóa sự khác biệt giữa bằng chứng nội bộ và bằng chứng láng giềng hay tham chiếu ngưỡng nền từ tập phát triển. Đồng thời, trong thực nghiệm AIOps, nhiều thiết kế phức tạp được gộp chung vào một hệ thống mà không tách kiểm định từng thành phần trên tập kiểm tra giữ kín, dễ che giấu những thành phần không mang lại lợi ích thực tế.
  \item \textbf{Khoảng trống G4 — Sự đứt gãy giữa thẩm định cảnh báo đã kiểm soát FP với phân tích nguyên nhân gốc có trích dẫn bằng chứng trên một nền tảng mở.} Phần lớn công trình hoặc dừng ở bước phát hiện/thẩm định cảnh báo, hoặc bắt đầu RCA từ danh sách sự cố đã xác nhận sẵn, và hiếm khi kết hợp cả hai trên một nền tảng mã nguồn mở có bộ nhớ đệm lời gọi LLM để tái lập thực nghiệm và trình diễn trực quan.
\end{enumerate}

Để lấp đầy bốn khoảng trống trên, đề tài đề xuất \textbf{\CTAV{}} (\emph{Calibrated Topology-Aware Alert Validation}) và tích hợp \CTAV{} làm tầng nền cho \textbf{tác tử phân tích nguyên nhân gốc} trên một nền tảng thống nhất. Cụ thể, trong không gian thiết kế của \CTAV{}, đề tài xây dựng và kiểm định tách biệt từng thành phần theo quy trình đăng ký giả thuyết trước: thành phần \textbf{C1} (truy xuất lai BM25--topology, đang dùng trong cấu hình cuối) kết hợp với thành phần \textbf{C3} (neo bằng chứng theo ngưỡng nền đo trên tập phát triển \splitname{dev}, đang dùng trong cấu hình cuối) đóng vai trò lõi giúp duy trì lợi ích nhận diện của topology đồng thời kìm mạnh báo động giả. Hai cơ chế bổ sung trong thiết kế ban đầu là \textbf{C2} (ngưỡng quyết định bất đối xứng theo nguồn bằng chứng, đã loại sau kiểm định) và \textbf{C4} (định tuyến hai tầng thích nghi, đã loại sau kiểm định) được đánh giá minh bạch trên tập \splitname{val} và loại khỏi cấu hình cuối vì không tạo lợi ích bổ sung so với cặp C1+C3 tinh gọn; tương tự, cơ chế \textbf{C5} (\code{Rescue Score}, chỉ đánh giá ở vai trò phân tích loại trừ) được khảo sát độc lập thay vì mặc định bật trong hệ thống chính. Phía sau \CTAV{}, tác tử RCA gom các cảnh báo đã thẩm định thành sự cố, lần ngược đường lan truyền trên đồ thị phụ thuộc dịch vụ và sinh báo cáo chẩn đoán có cấu trúc kèm trích dẫn bằng chứng. Toàn bộ thiết kế kiến trúc, thuật toán chi tiết và mô hình dữ liệu của hệ thống được trình bày trong Chương~\ref{ch:de-xuat}.

%        ...
%        \bibliographystyle{IEEEtran}
%        \bibliography{refs}
%
% Cac nhan lien chuong duoc tham chieu tu Chuong 3:
%   Chuong: ch:mo-dau, ch:nen-tang, ch:de-xuat, ch:thuc-nghiem, ch:platform
%   Chuong 2: sec:ch2-du-lieu, sec:ch2-truy-xuat, sec:ch2-judge, sec:ch2-rca,
%             sec:ch2-dang-ky
%   Chuong 4: sec:ch4-ctav, sec:ch4-c1, sec:ch4-policy, sec:ch4-c4, sec:ch4-rca
%   Chuong 5: sec:ch5-hypotheses, sec:ch5-c5, sec:ch5-budget
% =============================================================================

\chapter{Công trình liên quan}
\label{ch:cong-trinh-lien-quan}

Chương này hệ thống hóa các hướng nghiên cứu và các hệ thống vận hành thực tế liên quan trực tiếp tới bài toán thẩm định cảnh báo có nhận biết topology và phân tích nguyên nhân gốc trên kiến trúc microservice. Nội dung được tổ chức xoay quanh hai câu hỏi: cộng đồng học thuật và công nghiệp đã giải quyết những khía cạnh nào của chuỗi xử lý từ phát hiện bất thường đến định vị nguyên nhân gốc, và còn những khoảng trống nào mà đề xuất của đề tài nhắm tới. Mục~\ref{sec:ch3-nghien-cuu} khảo sát bảy nhánh nghiên cứu học thuật từ phát hiện bất thường chuỗi thời gian, ứng dụng mô hình ngôn ngữ lớn trong phân tích log và nguyên nhân gốc, kiến trúc giám sát lai, phân tích nguyên nhân gốc dựa trên đồ thị, truy xuất dựa trên đồ thị, hiệu chỉnh độ tin cậy của \emph{LLM-as-a-Judge}, đến gom nhóm và phân loại cảnh báo. Mục~\ref{sec:ch3-he-thong} đối chiếu các nền tảng AIOps thương mại và công cụ mã nguồn mở hiện hành dựa trên tài liệu chính thức. Cuối cùng, Mục~\ref{sec:ch3-ket-chuong} tổng hợp các hướng tiếp cận trong Bảng~\ref{tab:ch3-so-sanh}, rút ra bốn khoảng trống nghiên cứu và định vị rõ phạm vi đóng góp của \CTAV{} cùng tác tử phân tích nguyên nhân gốc.

% -----------------------------------------------------------------------------
\section{Các nghiên cứu liên quan}
\label{sec:ch3-nghien-cuu}

Các công trình liên quan đến đề tài nằm ở giao điểm giữa giám sát chuỗi thời gian, phân tích dữ liệu quan sát đa phương thức, truy xuất bằng chứng theo đồ thị và hiệu chỉnh mô hình ngôn ngữ lớn đóng vai trò bộ thẩm định. Bảy tiểu mục dưới đây phân tích từng nhánh theo cùng một trục: cách tiếp cận của công trình, hạn chế đặt trong bài toán của đề tài và điểm khác biệt của đề tài.

\subsection{Phát hiện bất thường trên chuỗi thời gian metric}
\label{sec:ch3-tsad}

Hướng phát hiện bất thường trên chuỗi thời gian không giám sát hình thành nền tảng cho tầng lọc cảnh báo đầu vào. Liu và cộng sự~\cite{liu2008iforest} đề xuất Isolation Forest, một phương pháp dựa trên tập cây phân tách ngẫu nhiên tận dụng nhận xét rằng điểm bất thường ít và khác biệt nên cần ít phép chia để cô lập hơn điểm bình thường; mô hình chạy nhanh, không cần nhãn và tiêu thụ ít tài nguyên. Mở rộng sang dữ liệu luồng có phân bố biến đổi theo thời gian, Guha và cộng sự~\cite{guha2016rrcf} xây dựng Robust Random Cut Forest (RRCF) cho phép cập nhật cấu trúc cây trực tuyến khi điểm mới đến và đo độ bất thường bằng mức thay đổi độ phức tạp cây khi chèn hoặc xóa điểm. Hạn chế chung của nhóm cây phân tách là chỉ nhìn từng vector metric độc lập theo cửa sổ, nên không phân biệt được dao động tải bình thường với sự cố thực sự nếu không đối chiếu ngữ cảnh văn bản từ log hoặc quan hệ gọi giữa các service.

Để nắm bắt tương quan đa biến và phụ thuộc giữa các cảm biến hay metric, các công trình học sâu dựa trên đồ thị đã được đề xuất. Zhao và cộng sự~\cite{zhao2020mtadgat} giới thiệu MTAD-GAT với hai lớp chú ý đồ thị song song: một lớp học tương quan giữa các đặc trưng tại cùng thời điểm và một lớp học phụ thuộc dọc trục thời gian, kết hợp đồng thời mục tiêu dự báo và tái tạo để tính điểm bất thường. Deng và Hooi~\cite{deng2021gdn} đề xuất Graph Deviation Network (GDN), tự học đồ thị quan hệ giữa các cảm biến qua vector nhúng và dùng mạng chú ý đồ thị để dự báo giá trị tiếp theo, từ đó giải thích bất thường qua độ lệch của các cạnh dự báo. Tuy đạt độ chính xác cao hơn trên nhiều tập chuỗi thời gian đa biến, nhóm học sâu đồ thị đòi hỏi huấn luyện lại khi topology hoặc tập metric thay đổi, cần tài nguyên tính toán lớn hơn đáng kể và vẫn chỉ khai thác tín hiệu số mà bỏ qua nội dung log hay trace vốn chứa thông tin tường minh về lỗi ứng dụng.

Khác với các công trình tập trung thay thế bộ phát hiện metric bằng kiến trúc học sâu ngày càng phức tạp, đề tài này giữ tầng phát hiện nhẹ bằng Isolation Forest để bảo đảm tốc độ và khả năng chạy lại tức thì trên mọi cấu hình, rồi chuyển trọng tâm sang tầng thẩm định bằng bằng chứng đa phương thức phía sau. Cách phân công này cho phép tận dụng tốc độ lọc sơ bộ của detector nhẹ trong khi vẫn kiểm soát được báo động giả nhờ đối chiếu log và đồ thị phụ thuộc service.

\subsection{Mô hình ngôn ngữ lớn trong phân tích log và nguyên nhân gốc}
\label{sec:ch3-llm-log-rca}

Sự xuất hiện của các mô hình ngôn ngữ lớn (LLM) mở ra hướng xử lý trực tiếp dữ liệu log phi cấu trúc và suy luận sự cố bằng ngôn ngữ tự nhiên. Ở bài toán phân tách mẫu log (\emph{log parsing}), Zhong và cộng sự~\cite{zhong2024logparser} đề xuất LogParser-LLM, kết hợp cây tiền tố với truy vấn LLM chọn lọc để trích xuất mẫu log mà không cần huấn luyện lại mô hình khi định dạng log thay đổi. Ở bài toán phát hiện bất thường trên chuỗi log, Guan và cộng sự~\cite{guan2024logllm} giới thiệu LogLLM, phối hợp bộ mã hóa ngữ nghĩa dạng BERT cho từng dòng log với bộ giải mã LLM nhằm mô hình hóa chuỗi log mà không cần bước tách mẫu thủ công. Hai công trình này giải quyết tốt khâu biểu diễn và phân tích log đơn lẻ, nhưng chỉ làm việc trên tín hiệu log tách rời mà không tích hợp luồng cảnh báo từ metric cũng như cấu trúc phụ thuộc giữa các service.

Ở bài toán phân tích nguyên nhân gốc cho sự cố đám mây, nhiều công trình sử dụng LLM trực tiếp trên hồ sơ sự cố hoặc dưới dạng tác tử gọi công cụ. Ahmed và cộng sự~\cite{ahmed2023recommending} thực hiện đánh giá quy mô lớn trên hơn 40.000 sự cố vận hành thực tế tại Microsoft để khuyến nghị nguyên nhân gốc và bước khắc phục bằng mô hình dòng GPT, cho thấy LLM có ích trong việc tóm tắt tri thức sự cố nhưng nhạy cảm với chất lượng mô tả đầu vào. Chen và cộng sự~\cite{chen2024rcacopilot} xây dựng RCACopilot, một hệ thống định tuyến sự cố vào quy trình thu thập dữ liệu chẩn đoán tương ứng rồi để LLM dự đoán nhóm nguyên nhân gốc kèm giải thích bằng ngôn ngữ tự nhiên. Tiến thêm bước tự chủ, Wang và cộng sự~\cite{wang2024rcagent} đề xuất RCAgent sử dụng tác tử LLM có hỗ trợ công cụ truy vấn log và tổng hợp tự nhất quán (\emph{self-consistency}) trong môi trường đám mây, trong khi Roy và cộng sự~\cite{roy2024exploring} khảo sát tác tử theo vòng lặp ReAct trang bị công cụ truy xuất log, metric và cơ sở dữ liệu lịch sử để điều tra sự cố.

Điểm hạn chế chung của nhóm công trình LLM cho RCA nêu trên là phần lớn giả định sự cố đầu vào đã được xác nhận hoặc đã mở phiếu sự cố sẵn, nên không đánh giá trực tiếp bài toán thẩm định luồng cảnh báo thô để chặn báo động giả trước khi gọi tác tử chẩn đoán tốn kém. Ngoài ra, các hệ thống đánh giá trên dữ liệu sự cố nội bộ doanh nghiệp khó tái lập độc lập bên ngoài. Đề tài này đặt tầng thẩm định cảnh báo có kiểm soát ngân sách ngay trước tác tử RCA và thực nghiệm hoàn toàn trên bộ dữ liệu công khai có thể tái lập.

\subsection{Kiến trúc giám sát lai và LLM-as-a-Judge}
\label{sec:ch3-hybrid-judge}

Để dung hòa chi phí gọi LLM với yêu cầu giám sát liên tục, nhiều công trình kết hợp mô hình nhẹ hoặc cơ chế truy xuất tăng cường sinh (\emph{Retrieval-Augmented Generation}, RAG)~\cite{lewis2020rag} với bộ phân xử LLM ở tầng sau. Về nguyên lý phân tầng chi phí, nền tảng lý thuyết bắt nguồn từ cơ chế đánh đổi sai số và quyền từ chối (\emph{reject option}) của Chow~\cite{chow1970reject} cùng mô hình chuỗi tầng gọi LLM tiết kiệm chi phí FrugalGPT của Chen và cộng sự~\cite{chen2024frugalgpt}, trong đó các mẫu rõ ràng được xử lý bởi mô hình rẻ hơn và chỉ vùng bất định mới chuyển lên mô hình đắt hơn.

Trong miền phát hiện bất thường, một loạt công trình gần đây đưa kiến trúc lai vào thực nghiệm. Russell-Gilbert và cộng sự~\cite{russellgilbert2025raad} đề xuất RAAD-LLM, tích hợp RAG vào LLM để thích nghi phát hiện bất thường chuỗi thời gian khi bối cảnh hệ thống thay đổi mà không cần tinh chỉnh trọng số. Gu và cộng sự~\cite{gu2025argos} giới thiệu ARGOS, một khuôn khổ tác tử cho phép LLM tự sinh và tinh chỉnh luật phát hiện bất thường trên chuỗi thời gian nhằm giảm phụ thuộc vào truy vấn trực tiếp trên từng mẫu. Devireddy và Huang~\cite{devireddy2025calm} đề xuất CALM, điều phối liên tục giữa phát hiện thời gian thực và cơ chế can thiệp thích nghi của LLM. Ở góc độ đánh giá năng lực thẩm định, Sivalingam và cộng sự~\cite{sivalingam2026tsquery} xây dựng TSQueryBench để kiểm tra khả năng dùng \emph{LLM-as-a-Judge} đánh giá lời giải thích bất thường chuỗi thời gian.

Hạn chế chung của các kiến trúc lai này khi áp vào microservice là chúng chủ yếu vận hành trên chuỗi thời gian đơn lẻ hoặc truy xuất ngữ cảnh văn bản phẳng, chưa kết nối tường minh với đồ thị phụ thuộc giữa các service sinh từ trace. Quan trọng hơn, các nghiên cứu này chưa khảo sát hệ quả ngược chiều khi mở rộng bằng chứng cho bộ thẩm định LLM: việc cấp thêm ngữ cảnh từ láng giềng có thể làm tăng tỷ lệ báo động giả do thiên lệch xác nhận. Đề tài này tập trung giải quyết trực tiếp bài toán kiểm soát ngân sách và kìm báo động giả trong kiến trúc lai detector nhẹ kết hợp \emph{LLM-as-a-Judge} cho hệ microservice.

\subsection{Phân tích nguyên nhân gốc dựa trên topology, đồ thị và suy luận nhân quả}
\label{sec:ch3-rca-graph}

Định vị nguyên nhân gốc dựa trên đồ thị phụ thuộc và suy luận nhân quả là nhánh phát triển lâu đời nhất trong giám sát hệ phân tán. Kim và cộng sự~\cite{kim2013monitorrank} đề xuất MonitorRank, thực hiện bước đi ngẫu nhiên (\emph{random walk}) có trọng số tương quan metric trên đồ thị gọi dịch vụ để xếp hạng service nghi vấn trong kiến trúc hướng dịch vụ. Chen và cộng sự~\cite{chen2014causeinfer} xây dựng CauseInfer với đồ thị nhân quả hai tầng (tầng dịch vụ và tầng metric) bằng thuật toán kiểm định độc lập có điều kiện PC để lần ngược nút gốc gây suy giảm hiệu năng. Trên kiến trúc microservice, Wu và cộng sự~\cite{wu2020microrca} giới thiệu MicroRCA, gắn dữ liệu sử dụng tài nguyên máy chủ vào đồ thị dịch vụ bất thường và áp dụng Personalized PageRank để định vị service gây lỗi hiệu năng.

Các công trình gần đây tiếp cận bài toán bằng mô hình nhân quả chặt chẽ hơn và dữ liệu đa phương thức. Li và cộng sự~\cite{li2022circa} đề xuất CIRCA, mô hình hóa lỗi như một can thiệp (\emph{intervention}) trên mạng Bayes nhân quả xây từ kiến trúc hệ thống và dùng điều kiện thay đổi phân bố có điều kiện để chỉ ra chỉ số gốc. Ikram và cộng sự~\cite{ikram2022rcd} đề xuất RCD, xem sự cố như can thiệp cục bộ và khám phá đồ thị nhân quả theo phân cấp để tránh chi phí tính toán lớn của thuật toán PC trên toàn bộ hệ thống. Kết hợp cả ba trụ cột dữ liệu quan sát, Yu và cộng sự~\cite{yu2023nezha} đề xuất Nezha, chuyển đổi metric, log và trace thành biểu diễn sự kiện đồng nhất và so sánh mẫu phân kỳ giữa trạng thái bình thường với trạng thái lỗi để định vị nguyên nhân gốc ở mức chi tiết, đồng thời công bố bộ dữ liệu microservice đa phương thức làm chuẩn đánh giá. Đáng lưu ý về tính bền vững của chuỗi xử lý, Pham và cộng sự~\cite{pham2024baro} đề xuất BARO bằng phát hiện điểm đổi Bayes đa biến kết hợp kiểm định phi tham số, đồng thời chỉ ra rằng sai số ở khâu phát hiện bất thường đầu vào có thể lan sang làm giảm mạnh độ chính xác định vị gốc; tiếp đó nhóm tác giả công bố bộ chuẩn RCAEval~\cite{pham2025rcaeval} tổng hợp 735 ca lỗi đa phương thức trên nhiều hệ microservice mã nguồn mở để thống nhất việc so sánh.

Hạn chế của nhóm phương pháp đồ thị và nhân quả thuần túy gồm hai điểm. Thứ nhất, phần lớn mô hình chỉ xuất ra danh sách xếp hạng service hoặc metric gốc dưới dạng chỉ số toán học, thiếu báo cáo tường minh giải thích cơ chế lan truyền bằng chứng từ log và trace cho kỹ sư vận hành. Thứ hai, trừ những phân tích gần đây như BARO, đa số nghiên cứu RCA giả định cửa sổ thời gian sự cố đã được xác định sạch từ trước, tách rời khỏi bài toán lọc cảnh báo giả trực tuyến. Đề tài này nối trực tiếp hai khâu: dùng \CTAV{} để làm sạch luồng cảnh báo trước khi gom sự cố, rồi dùng tác tử RCA lần ngược chuỗi lan truyền trên đồ thị phụ thuộc và sinh báo cáo có cấu trúc kèm trích dẫn bằng chứng.

\subsection{Truy xuất dựa trên đồ thị và GraphRAG cho vận hành hệ thống}
\label{sec:ch3-graphrag}

Trong lĩnh vực xử lý ngôn ngữ tự nhiên, Edge và cộng sự~\cite{edge2024graphrag} đề xuất GraphRAG nhằm khắc phục hạn chế của truy xuất văn bản phẳng khi cần tổng hợp thông tin phân tán trên nhiều tài liệu. Phương pháp xây dựng đồ thị tri thức thực thể--quan hệ từ kho văn bản, phân cụm cộng đồng đồ thị và sinh tóm tắt phân cấp để trả lời các truy vấn mang tính toàn cục. Tuy mở ra hướng truy xuất có cấu trúc, GraphRAG thiết kế cho kho tài liệu tĩnh với đồ thị trích xuất bằng LLM, nên vừa tốn chi phí lập chỉ mục vừa khác về bản chất so với đồ thị phụ thuộc động của hệ microservice được đo trực tiếp từ quan hệ gọi \code{SpanID}--\code{ParentID} trong trace.

Khi áp dụng tư tưởng truy xuất theo cấu trúc đồ thị vào bài toán thẩm định cảnh báo microservice, một nghiên cứu tiền đề về truy xuất bằng chứng có nhận biết topology~\cite{author2026icaetm} đã đề xuất kết hợp điểm từ khóa BM25 với độ gần về thời gian và khoảng cách bước nhảy trên đồ thị phụ thuộc service để cung cấp thêm bằng chứng từ các service láng giềng cho bộ thẩm định LLM (chi tiết công thức hợp nhất đã trình bày ở Mục~\ref{sec:ch2-truy-xuat}). Kết quả định tính của hướng tiếp cận này cho thấy một đánh đổi rõ rệt: việc mở rộng ngữ cảnh sang các service có liên kết gọi giúp cải thiện khả năng nhận diện các sự cố lan truyền mà bằng chứng cục bộ không đủ, nhưng khi đưa trực tiếp vào \emph{LLM-as-a-Judge} mà không có cơ chế kiểm soát chính sách, số báo động giả tăng lên mạnh do bộ thẩm định dễ diễn giải nhầm log nền từ láng giềng thành dấu hiệu bất thường. Đề tài này xuất phát trực tiếp từ đánh đổi đó để xây dựng cơ chế hiệu chỉnh theo nguồn gốc bằng chứng và kiểm định nghiêm ngặt hiệu quả kìm báo động giả trên dữ liệu ngoài tập tinh chỉnh.

\subsection{Hiệu chỉnh độ tin cậy, thiên lệch của LLM-as-a-Judge và gắn bằng chứng}
\label{sec:ch3-calibration}

Dù \emph{LLM-as-a-Judge} đạt mức đồng thuận cao với người đánh giá trên nhiều tác vụ tổng quát~\cite{zheng2023judging}, hàng loạt nghiên cứu đã chứng minh mô hình ngôn ngữ mắc nhiều thiên lệch có hệ thống và thường tự tin thái quá. Wang và cộng sự~\cite{wang2024fair} chỉ ra rằng bộ thẩm định LLM chịu thiên lệch vị trí nghiêm trọng khi so sánh cặp câu trả lời, dẫn đến thứ hạng dễ bị thao túng chỉ bằng cách đảo thứ tự trình bày, và đề xuất khung FairEval hiệu chỉnh bằng cách buộc sinh bằng chứng trước khi chấm điểm kết hợp trung bình hóa vị trí. Sharma và cộng sự~\cite{sharma2024sycophancy} phân tích hiện tượng chiều lòng (\emph{sycophancy}), cho thấy mô hình đã qua tinh chỉnh theo phản hồi con người có xu hướng đồng thuận với giả định ngầm trong câu hỏi ngay cả khi giả định đó sai — một cơ chế tương đồng trực tiếp với thiên lệch xác nhận khi bộ thẩm định nhận thông báo rằng ``cảnh báo đang kích hoạt''.

Về hiệu chỉnh xác suất và ước lượng độ bất định, Guo và cộng sự~\cite{guo2017calibration} chứng minh các mạng nơ-ron hiện đại tuy chính xác nhưng hiệu chỉnh kém và đề xuất phương pháp hiệu chỉnh hậu xử lý bằng tỉ lệ nhiệt độ (\emph{temperature scaling}). Trên LLM, Kadavath và cộng sự~\cite{kadavath2022mostly} khảo sát khả năng mô hình tự đánh giá xác suất đúng của phát biểu do chính nó sinh ra, cho thấy tín hiệu nội tại có thể hiệu chỉnh được trong câu hỏi trắc nghiệm nhưng khó khái quát khi chuyển phân bố tác vụ. Gần hơn với bài toán sử dụng API thương mại không truy cập được xác suất token thô, Xiong và cộng sự~\cite{xiong2024uncertainty} đánh giá thực nghiệm độ tin cậy bằng lời (\emph{verbalized confidence}) trên nhiều dòng LLM và chỉ ra mô hình thường gán điểm tin cậy quá cao, tập trung dày đặc ở vùng sát $1{,}0$. Song song đó, để hạn chế suy diễn vô căn cứ trong văn bản sinh ra, Gao và cộng sự~\cite{gao2023alce} đề xuất chuẩn đánh giá ALCE đo đồng thời tính đúng đắn và chất lượng trích dẫn nguồn bằng chứng trong câu trả lời của LLM.

Khoảng trống còn lại trong nhánh nghiên cứu này là các kỹ thuật hiệu chỉnh hiện có chủ yếu thiết kế cho bài toán chấm điểm cặp câu trả lời, hỏi--đáp mở hoặc sinh văn bản có trích dẫn tổng quát. Chưa có cơ chế nào hiệu chỉnh trực tiếp quyết định của \emph{LLM-as-a-Judge} trong bài toán thẩm định cảnh báo microservice theo \emph{nguồn gốc bằng chứng} (nội bộ hay láng giềng trên đồ thị phụ thuộc) và theo \emph{mức độ bất thường nền đo trên tập phát triển}. Đề tài này cụ thể hóa ý tưởng hiệu chỉnh bằng chính sách neo bằng chứng và kiểm định tách biệt từng thành phần thiết kế (ngưỡng bất đối xứng theo nguồn bằng chứng, neo ngưỡng nền trên tập phát triển và định tuyến hai tầng) theo quy trình đăng ký giả thuyết trước ở Mục~\ref{sec:ch2-dang-ky}.

\subsection{Gom nhóm, phân loại và cứu cảnh báo trong vận hành}
\label{sec:ch3-alert-triage}

Trong vận hành hệ thống trực tuyến quy mô lớn, xử lý bão cảnh báo và định tuyến sự cố là bài toán sống còn để giảm tải cho kỹ sư trực. Zhao và cộng sự~\cite{zhao2020alertstorm} khảo sát hiện tượng bão cảnh báo trong hệ dịch vụ trực tuyến công nghiệp và đề xuất quy trình hai bước gồm phát hiện thời điểm bão và gom nhóm các cảnh báo tương quan về thời gian--ngữ nghĩa thành một cụm sự cố đại diện. Ở khâu định tuyến sự cố, Chen và cộng sự~\cite{chen2019empirical,chen2019continuous} phân tích nguyên nhân định tuyến nhầm sự cố tại các dịch vụ đám mây lớn và xây dựng phương pháp phân loại liên tục tận dụng thông tin cập nhật trong quá trình thảo luận của kỹ sư. Gần đây hơn, Kuang và cộng sự~\cite{kuang2024knowledge} đề xuất phương pháp gom cảnh báo có nhận biết tri thức cho hệ đám mây quy mô lớn, kết hợp đặc trưng văn bản với tri thức vận hành và topology để gộp cảnh báo cùng gốc tốt hơn các luật tĩnh.

Các công trình gom nhóm và phân loại này giải quyết tốt khâu \emph{sau khi cảnh báo đã phát sinh}, giúp nén số lượng thông báo gửi tới người trực, nhưng thường không thẩm định lại bản thân cảnh báo đầu vào có phải là báo động giả do bộ phát hiện metric gây ra hay không. Mặt khác, một câu hỏi ít được kiểm chứng nghiêm ngặt là liệu các cửa sổ bị bộ phát hiện chính bỏ sót có thể được ``cứu'' lại bằng điểm phụ trợ từ tín hiệu log/trace (cơ chế \code{Rescue Score}) mà không làm vỡ tỷ lệ báo động giả hay không. Đề tài này tách rõ hai vấn đề: sử dụng gom nhóm theo cửa sổ và topology ở đầu vào tác tử RCA để hợp nhất các cảnh báo đã được \CTAV{} thẩm định thành sự cố, đồng thời đánh giá cơ chế cứu cảnh báo như một phân tích loại trừ độc lập (Mục~\ref{sec:ch5-c5}) thay vì mặc định đưa vào hệ thống chính.

% -----------------------------------------------------------------------------
\section{Các hệ thống và công cụ hiện có}
\label{sec:ch3-he-thong}

Bên cạnh các công trình học thuật, nhiều nền tảng AIOps thương mại và công cụ mã nguồn mở đã được triển khai rộng rãi trong thực tế vận hành. Mục này mô tả tính năng cốt lõi của các hệ thống tiêu biểu dựa trên tài liệu chính thức của nhà phát triển, sau đó chỉ ra điểm khác biệt về mục tiêu và giới hạn khi đặt trong bối cảnh nghiên cứu của đề tài.

\subsection{Nền tảng AIOps thương mại}
\label{sec:ch3-thuong-mai}

\textbf{Datadog (Event Management và Bits AI).} Theo tài liệu chính thức của Datadog~\cite{datadog2026docs}, phân hệ Event Management thu nạp tập trung cảnh báo từ các bộ giám sát nội bộ cùng nguồn bên thứ ba, sau đó áp dụng tương quan cảnh báo tự động dựa trên topology dịch vụ, nhãn chia sẻ và mẫu thời gian để gom các cảnh báo lan truyền vào chung một sự vụ, đồng thời làm giàu ngữ cảnh về đơn vị sở hữu dịch vụ. Ở lớp điều tra sự cố, tác tử Bits AI (Bits Investigation) được kích hoạt tự động khi có cảnh báo nghiêm trọng, thực hiện kiểm định song song nhiều giả thuyết trên dữ liệu metric, log, trace và thay đổi cấu hình để lần theo chuỗi nhân quả và đề xuất hướng xử lý. Tuy cung cấp quy trình khép kín từ gom cảnh báo đến điều tra bằng AI, kiến trúc bên trong của cơ chế tương quan và điều tra đều đóng kín; nhà cung cấp không công bố phương pháp hiệu chỉnh độ tin cậy của bộ suy luận LLM theo nguồn bằng chứng hay đánh giá tách biệt tác động của topology lên tỷ lệ báo động giả trên bộ chuẩn công khai.

\textbf{Dynatrace (Davis AI và Smartscape).} Tài liệu kỹ thuật của Dynatrace~\cite{dynatrace2026docs} mô tả kiến trúc dựa trên hai trụ cột là bản đồ phụ thuộc thời gian thực Smartscape và động cơ phân tích nhân quả Davis AI. Smartscape tự động lập đồ thị phụ thuộc đầy đủ từ ứng dụng, service, tiến trình đến máy chủ và hạ tầng Kubernetes. Khi phát sinh bất thường ở mức chỉ số, độ trễ hoặc tỷ lệ lỗi, Davis AI gom các sự kiện liên quan có cùng nguyên nhân gốc thành một thực thể \emph{Problem} duy nhất và thực hiện phân tích cây lỗi tất định (\emph{deterministic fault-tree analysis}) trên đồ thị Smartscape để tách biệt thành phần gốc khỏi các thành phần chịu ảnh hưởng, hiển thị qua biểu đồ \emph{Visual Resolution Path}. Ưu điểm của Dynatrace là tính nhất quán cao nhờ mô hình nhân quả dựa trên agent thu thập sâu (\emph{OneAgent}), nhưng hệ thống gắn chặt với hạ tầng thu thập độc quyền, chi phí bản quyền lớn và không nhằm giải bài toán hiệu chỉnh quyết định của \emph{LLM-as-a-Judge} trên dữ liệu log phi cấu trúc lẫn ngân sách truy vấn LLM giới hạn.

\textbf{PagerDuty AIOps.} Theo tài liệu hỗ trợ chính thức của PagerDuty~\cite{pagerduty2026docs}, bộ công cụ PagerDuty AIOps tập trung vào ba chức năng: \emph{Intelligent Alert Grouping} dùng học máy nhóm các cảnh báo có liên quan theo mẫu lịch sử để giảm nhiễu thông báo; \emph{Probable Origin} phân tích chuỗi sự cố lan truyền trong lịch sử sáu tháng gần nhất để gợi ý tối đa ba dịch vụ hoặc sự cố có khả năng là điểm khởi phát cao nhất; và \emph{Event Orchestration} cung cấp bộ máy luật cho phép lọc, làm giàu và định tuyến sự kiện tự động ở mức toàn cục hoặc từng dịch vụ. Cách tiếp cận của PagerDuty giải quyết hiệu quả bài toán quản trị phản ứng sự cố ở tầng thông báo, nhưng phụ thuộc nhiều vào lịch sử sự cố tích lũy theo tổ chức, không khai thác đồ thị gọi chi tiết từng cửa sổ trace để thẩm định cảnh báo từ tín hiệu thô và không công bố cơ chế kiểm soát báo động giả cho tác tử thẩm định LLM.

\subsection{Công cụ mã nguồn mở}
\label{sec:ch3-ma-nguon-mo}

\textbf{Prometheus Alertmanager.} Theo tài liệu chính thức của dự án Prometheus~\cite{prometheus2026alertmanager}, Alertmanager đảm nhận tầng xử lý thông báo sau khi các luật cảnh báo Prometheus kích hoạt, cung cấp bốn cơ chế nền tảng: cây định tuyến theo nhãn (\code{route}), gom nhóm cảnh báo cùng nhãn thành một thông báo (\code{group\_by}), ức chế cảnh báo đích khi cảnh báo nguồn cùng nhóm đang kích hoạt (\code{inhibit\_rules}) và tắt tiếng tạm thời theo khung bảo trì (\code{silences}). Alertmanager nhẹ, minh bạch và là tiêu chuẩn thực tế trong hệ sinh thái đám mây, nhưng mọi quan hệ ức chế và gom nhóm đều phải cấu hình tĩnh bằng tay trong tệp YAML; công cụ không tự suy diễn đồ thị phụ thuộc từ trace, không đọc nội dung log và không thực hiện thẩm định cảnh báo hay phân tích nguyên nhân gốc.

\textbf{Keep.} Keep là nền tảng AIOps và quản lý cảnh báo mã nguồn mở; theo tài liệu chính thức của dự án~\cite{keep2026docs}, hệ thống hỗ trợ khử trùng lặp cảnh báo bằng dấu vân tay (\emph{fingerprint deduplication}, gồm chế độ toàn phần và từng phần), cung cấp ba cơ chế tương quan cảnh báo (luật thủ công bằng biểu thức CEL, tương quan theo topology và tương quan hỗ trợ bởi AI) cùng quy trình tự động hóa khai báo bằng YAML để làm giàu dữ liệu và gọi công cụ bên ngoài. Keep khắc phục được tính đơn điệu của công cụ luật tĩnh và mở rộng kết nối tốt với nhiều nguồn giám sát, nhưng vẫn xoay quanh bài toán tổng hợp--điều phối cảnh báo và quy trình tự động hóa; tài liệu hiện hành chưa đề cập cơ chế hiệu chỉnh \emph{LLM-as-a-Judge} theo nguồn bằng chứng hay khung thực nghiệm tái lập đánh giá đánh đổi recall--FPR dưới ngân sách LLM cố định.

\textbf{HolmesGPT.} HolmesGPT là tác tử điều tra sự cố mã nguồn mở khởi nguồn từ Robusta và hiện là dự án CNCF Sandbox; theo tài liệu chính thức~\cite{holmesgpt2026docs}, công cụ vận hành theo vòng lặp ReAct chạy ngay trong cụm Kubernetes ở chế độ chỉ đọc và tuân thủ phân quyền RBAC. Khi nhận cảnh báo từ Alertmanager hoặc PagerDuty, HolmesGPT để LLM tự chọn và gọi các công cụ truy vấn dữ liệu quan sát (như Prometheus, Loki hoặc \code{kubectl}), đọc kết quả trả về qua nhiều bước rồi tổng hợp nhận định nguyên nhân gốc. HolmesGPT minh họa rõ tính khả thi của tác tử LLM mã nguồn mở trong vận hành thực tế, nhưng chủ yếu đảm nhận khâu điều tra sau khi cảnh báo đã gửi tới, tốn nhiều lượt gọi LLM cho mỗi cảnh báo theo vòng lặp mở, và chưa tích hợp tầng thẩm định hiệu chỉnh trước đó để chặn báo động giả hay khung đánh giá thực nghiệm thống kê trên bộ chuẩn công khai.

\subsection{Nhận xét về các hệ thống hiện có}
\label{sec:ch3-nhan-xet}

Qua đối chiếu ba nền tảng thương mại và ba công cụ mã nguồn mở, có thể rút ra hai nhận xét chính về hiện trạng công nghiệp. Thứ nhất, các nền tảng thương mại (Datadog, Dynatrace, PagerDuty) đã tích hợp sâu khả năng gom cảnh báo theo topology và từng bước đưa tác tử AI vào điều tra sự cố, nhưng tất cả đều là hệ thống hộp đen gắn với hệ sinh thái thu thập độc quyền; chúng không cho phép kiểm định độc lập tác động của từng thành phần truy xuất hay chính sách hiệu chỉnh lên tỷ lệ báo động giả và không phục vụ được mục tiêu tái lập thực nghiệm khoa học. Thứ hai, ở mảng mã nguồn mở, hai năng lực cốt lõi vẫn nằm tách rời: nhóm quản lý cảnh báo (Alertmanager, Keep) mạnh về khử trùng lặp, luật ức chế và quy trình gom nhóm nhưng thiếu bộ thẩm định dựa trên bằng chứng log--trace đã hiệu chỉnh, còn tác tử điều tra bằng LLM (HolmesGPT) lại chạy trực tiếp theo vòng lặp nhiều bước mà không có tầng lọc kiểm soát ngân sách và kìm báo động giả ở đầu vào.

Do đó, thay vì cạnh tranh về độ phủ kết nối hạ tầng sản xuất với các sản phẩm thương mại, nền tảng trong đề tài này được xây dựng như một kiến trúc mở gọn nhẹ phục vụ đồng thời hai mục đích: tái lập chặt chẽ các thí nghiệm so sánh và loại trừ trên bộ dữ liệu công khai, và trình diễn trọn vẹn luồng xử lý từ phát lại dữ liệu quan sát, thẩm định cảnh báo bằng \CTAV{} đến định vị nguyên nhân gốc có bằng chứng.

% -----------------------------------------------------------------------------
\section{Kết chương}
\label{sec:ch3-ket-chuong}

Mục này tổng hợp các nghiên cứu học thuật ở Mục~\ref{sec:ch3-nghien-cuu} và các hệ thống thực tế ở Mục~\ref{sec:ch3-he-thong} trên cùng một hệ tiêu chí, qua đó chỉ rõ bốn khoảng trống nghiên cứu và định vị đóng góp của đề tài.

\subsection{Bảng so sánh tổng hợp}
\label{sec:ch3-bang-so-sanh}

Bảng~\ref{tab:ch3-so-sanh} đối chiếu các nhóm công trình tiêu biểu và các hệ thống hiện có theo năm tiêu chí bám sát trực tiếp mục tiêu của đề tài: (1) khai thác đồ thị topology dịch vụ, (2) sử dụng LLM để thẩm định hoặc suy luận sự cố, (3) hiệu chỉnh quyết định theo nguồn gốc bằng chứng, (4) thực hiện phân tích nguyên nhân gốc (RCA), và (5) có đánh giá định lượng báo động giả (FP/FPR) trên quy trình mở tái lập được. Trong bảng, \emph{Có} biểu thị tiêu chí được đáp ứng trực tiếp; \emph{Một phần} biểu thị tiêu chí chỉ được đáp ứng gián tiếp hoặc ở phạm vi hẹp (chẳng hạn chỉ có luật ức chế tĩnh, chỉ học đồ thị cảm biến không phải topology dịch vụ, hoặc hệ thống thương mại đóng nguồn không công bố kiểm định FP độc lập); \emph{Không} biểu thị tiêu chí nằm ngoài phạm vi của công trình hoặc công cụ.

\begin{table}[htbp]
  \centering
  \caption[So sánh công trình liên quan, hệ thống hiện có và đề xuất]{So sánh các nhóm nghiên cứu liên quan, hệ thống hiện có và đề xuất của đề tài.}
  \label{tab:ch3-so-sanh}
  \footnotesize
  \renewcommand{\arraystretch}{1.15}
  \setlength{\tabcolsep}{3.5pt}
  \begin{tabularx}{\linewidth}{@{}Y >{\centering\arraybackslash}p{1.45cm} >{\centering\arraybackslash}p{1.45cm} >{\centering\arraybackslash}p{1.6cm} >{\centering\arraybackslash}p{1.2cm} >{\centering\arraybackslash}p{1.65cm}@{}}
    \toprule
    \textbf{Công trình / Hệ thống} &
    \textbf{Topology} &
    \textbf{LLM judge} &
    \textbf{Hiệu chỉnh BC} &
    \textbf{RCA} &
    \textbf{Đánh giá FP \& tái lập} \\
    \midrule
    Isolation Forest~\cite{liu2008iforest}, RRCF~\cite{guha2016rrcf} &
      Không & Không & Không & Không & Một phần \\
    MTAD-GAT~\cite{zhao2020mtadgat}, GDN~\cite{deng2021gdn} &
      Một phần & Không & Không & Không & Một phần \\
    LogParser-LLM~\cite{zhong2024logparser}, LogLLM~\cite{guan2024logllm} &
      Không & Có & Không & Không & Một phần \\
    RCACopilot~\cite{chen2024rcacopilot}, RCAgent~\cite{wang2024rcagent}, Roy~\cite{roy2024exploring} &
      Không & Có & Không & Có & Không \\
    RAAD-LLM~\cite{russellgilbert2025raad}, ARGOS~\cite{gu2025argos}, CALM~\cite{devireddy2025calm} &
      Không & Có & Không & Không & Một phần \\
    MonitorRank~\cite{kim2013monitorrank}, MicroRCA~\cite{wu2020microrca}, CIRCA~\cite{li2022circa}, RCD~\cite{ikram2022rcd} &
      Có & Không & Không & Có & Không \\
    Nezha~\cite{yu2023nezha}, BARO~\cite{pham2024baro}, RCAEval~\cite{pham2025rcaeval} &
      Có & Không & Không & Có & Một phần \\
    GraphRAG~\cite{edge2024graphrag} &
      Một phần & Có & Không & Không & Không \\
    Truy xuất topology cho thẩm định LLM~\cite{author2026icaetm} &
      Có & Có & Không & Không & Có \\
    FairEval~\cite{wang2024fair}, Xiong~\cite{xiong2024uncertainty}, ALCE~\cite{gao2023alce} &
      Không & Có & Một phần & Không & Một phần \\
    Alert Storm~\cite{zhao2020alertstorm}, Kuang~\cite{kuang2024knowledge} &
      Một phần & Không & Không & Không & Một phần \\
    \midrule
    Datadog (Event Mgmt + Bits AI)~\cite{datadog2026docs} &
      Có & Có & Không & Có & Không \\
    Dynatrace (Davis AI + Smartscape)~\cite{dynatrace2026docs} &
      Có & Không & Không & Có & Không \\
    PagerDuty AIOps~\cite{pagerduty2026docs} &
      Một phần & Không & Không & Một phần & Không \\
    Prometheus Alertmanager~\cite{prometheus2026alertmanager} &
      Không & Không & Không & Không & Một phần \\
    Keep~\cite{keep2026docs} &
      Có & Một phần & Không & Không & Một phần \\
    HolmesGPT~\cite{holmesgpt2026docs} &
      Một phần & Có & Không & Có & Không \\
    \midrule
    \textbf{Đề tài này (\CTAV{} + tác tử RCA)} &
      \textbf{Có} & \textbf{Có} & \textbf{Có} & \textbf{Có} & \textbf{Có} \\
    \bottomrule
  \end{tabularx}
\end{table}

Bảng~\ref{tab:ch3-so-sanh} cho thấy bức tranh phân mảnh rõ rệt giữa các hướng hiện có. Nhóm phát hiện bất thường chuỗi thời gian và kiến trúc giám sát lai chưa khai thác cấu trúc đồ thị dịch vụ; nhóm RCA dựa trên đồ thị và nhân quả khai thác tốt topology nhưng không dùng LLM thẩm định luồng cảnh báo thô để kìm báo động giả; nhóm tác tử RCA dùng LLM bỏ qua khâu lọc cảnh báo giả có kiểm soát ngân sách; còn nhóm hiệu chỉnh \emph{LLM-as-a-Judge} mới xét bài toán ngôn ngữ tổng quát chứ chưa hiệu chỉnh theo nguồn gốc bằng chứng trên đồ thị phụ thuộc. Các hệ thống công nghiệp hoặc đóng nguồn không thể tái lập kiểm định FP, hoặc tách rời tầng quản lý cảnh báo với tầng điều tra LLM.

\subsection{Khoảng trống nghiên cứu và định vị đề xuất}
\label{sec:ch3-khoang-trong}

Từ kết quả khảo sát và Bảng~\ref{tab:ch3-so-sanh}, đề tài xác định bốn khoảng trống nghiên cứu cụ thể:

\begin{enumerate}
  \item \textbf{Khoảng trống G1 — Khoảng cách giữa phát hiện tốc độ cao và thẩm định chính xác dưới ngân sách LLM hữu hạn.} Bộ phát hiện metric không giám sát chạy nhanh nhưng dễ tạo báo động giả hoặc bỏ sót lỗi ngữ nghĩa; ngược lại, tác tử LLM gọi nhiều công cụ trực tiếp quá tốn kém để chạy trên mọi cửa sổ thời gian. Cần một cơ chế kết nối có kiểm soát chặt số lượt gọi LLM mà vẫn bảo đảm chất lượng thẩm định.
  \item \textbf{Khoảng trống G2 — Đánh đổi chưa được kiểm soát giữa thu hồi theo topology và báo động giả do thiên lệch xác nhận.} Mở rộng truy xuất bằng chứng sang láng giềng trên đồ thị phụ thuộc giúp nhận diện các lỗi lan truyền mà truy xuất từ khóa cục bộ bỏ lỡ, nhưng đồng thời khuếch đại xu hướng xác nhận nhầm của \emph{LLM-as-a-Judge} trên các cửa sổ bình thường nếu không có chính sách lọc và neo bằng chứng phù hợp.
  \item \textbf{Khoảng trống G3 — Thiếu cơ chế hiệu chỉnh bộ thẩm định cảnh báo theo nguồn gốc bằng chứng và thiếu kiểm chứng loại trừ có đăng ký giả thuyết trước.} Các phương pháp hiệu chỉnh LLM hiện hành chưa mô hình hóa sự khác biệt giữa bằng chứng nội bộ và bằng chứng láng giềng hay tham chiếu ngưỡng nền từ tập phát triển. Đồng thời, trong thực nghiệm AIOps, nhiều thiết kế phức tạp được gộp chung vào một hệ thống mà không tách kiểm định từng thành phần trên tập kiểm tra giữ kín, dễ che giấu những thành phần không mang lại lợi ích thực tế.
  \item \textbf{Khoảng trống G4 — Sự đứt gãy giữa thẩm định cảnh báo đã kiểm soát FP với phân tích nguyên nhân gốc có trích dẫn bằng chứng trên một nền tảng mở.} Phần lớn công trình hoặc dừng ở bước phát hiện/thẩm định cảnh báo, hoặc bắt đầu RCA từ danh sách sự cố đã xác nhận sẵn, và hiếm khi kết hợp cả hai trên một nền tảng mã nguồn mở có bộ nhớ đệm lời gọi LLM để tái lập thực nghiệm và trình diễn trực quan.
\end{enumerate}

Để lấp đầy bốn khoảng trống trên, đề tài đề xuất \textbf{\CTAV{}} (\emph{Calibrated Topology-Aware Alert Validation}) và tích hợp \CTAV{} làm tầng nền cho \textbf{tác tử phân tích nguyên nhân gốc} trên một nền tảng thống nhất. Cụ thể, trong không gian thiết kế của \CTAV{}, đề tài xây dựng và kiểm định tách biệt từng thành phần theo quy trình đăng ký giả thuyết trước: thành phần \textbf{C1} (truy xuất lai BM25--topology, đang dùng trong cấu hình cuối) kết hợp với thành phần \textbf{C3} (neo bằng chứng theo ngưỡng nền đo trên tập phát triển \splitname{dev}, đang dùng trong cấu hình cuối) đóng vai trò lõi giúp duy trì lợi ích nhận diện của topology đồng thời kìm mạnh báo động giả. Hai cơ chế bổ sung trong thiết kế ban đầu là \textbf{C2} (ngưỡng quyết định bất đối xứng theo nguồn bằng chứng, đã loại sau kiểm định) và \textbf{C4} (định tuyến hai tầng thích nghi, đã loại sau kiểm định) được đánh giá minh bạch trên tập \splitname{val} và loại khỏi cấu hình cuối vì không tạo lợi ích bổ sung so với cặp C1+C3 tinh gọn; tương tự, cơ chế \textbf{C5} (\code{Rescue Score}, chỉ đánh giá ở vai trò phân tích loại trừ) được khảo sát độc lập thay vì mặc định bật trong hệ thống chính. Phía sau \CTAV{}, tác tử RCA gom các cảnh báo đã thẩm định thành sự cố, lần ngược đường lan truyền trên đồ thị phụ thuộc dịch vụ và sinh báo cáo chẩn đoán có cấu trúc kèm trích dẫn bằng chứng. Toàn bộ thiết kế kiến trúc, thuật toán chi tiết và mô hình dữ liệu của hệ thống được trình bày trong Chương~\ref{ch:de-xuat}.
